% SPDX-License-Identifier: Apache-2.0
\documentclass[11pt]{article}
\setlength{\oddsidemargin}{0pt}
\setlength{\evensidemargin}{0pt}
\setlength{\textwidth}{6.5in}
\setlength{\textheight}{9in}
\setlength{\topmargin}{-0.5in}
\usepackage{amsmath,amssymb,hyperref}
% Use available Roman, bold and math-italic metrics in the built-in compiler.
\DeclareFontShape{OT1}{cmr}{m}{n}{<-> cmr10}{}
\DeclareFontShape{OT1}{cmr}{bx}{n}{<-> cmbx10}{}
\DeclareFontShape{OML}{cmm}{m}{it}{<-> cmmi10}{}
% Draw the two proof-end markers without the unavailable AMS square glyph.
\renewcommand{\square}{\mathord{\vcenter{\hrule\hbox{\vrule height .45em\kern .45em\vrule}\hrule}}}
\newcommand{\toprule}{\hline}
\newcommand{\midrule}{\hline}
\newcommand{\bottomrule}{\hline}
\hypersetup{hidelinks}
\setlength{\parindent}{0pt}
\setlength{\parskip}{0.55em}
\newcommand{\QQ}{\mathbb Q}
\newcommand{\ZZ}{\mathbb Z}
\newcommand{\GL}{\mathrm{GL}}
\newcommand{\Sp}{\mathrm{Sp}}
\newcommand{\Wh}{\mathrm{Wh}}
\newcommand{\N}{\mathrm N}
\newcommand{\eps}{\varepsilon}
\title{Metaplectic forms beyond the cubic QRH construction\\
\large A structural and analytic research map}
\author{Research review for the prime-factor oscillations project}
\date{4 November 2026}
\begin{document}
\maketitle
{\small Copyright (c) 2026 Prime Factor Oscillations contributors.
This document is released under the
\href{https://www.apache.org/licenses/LICENSE-2.0}{Apache License 2.0 (Apache-2.0)}.}
\begin{abstract}
We classify the choices that matter when seeking a metaplectic extension
of the QRH method: covering datum, group, field, automorphic representation,
Fourier model, coefficient support, transformation law and analytic consumer.
The cubic construction is a threefold theta representation used with sextic
characters. Its scalar coefficient identity is not a theorem for every
higher cover. We identify alternatives on general linear, symplectic and
exceptional groups, distinguish earlier project attempts from untested
cases, and give quantitative tests for a candidate's usefulness.

We compute the complete sevenfold and twelvefold $G_2$ and eightfold
$\Sp_4$ functional-coordinate models, retaining every character-grade
cross term. The sevenfold complete energy is independent of all Gauss
phases; a principal-preserving cubic correlation recovers reciprocal
sevenfold phase factors at the same torus bidegree. The eightfold case
has a unique positive normalized central mixture canceling its free
Jacobi phase with both period factors intact.
All eight sevenfold Jacobi channels admit individual compensated
readouts and normalized Abel limits. Their six-channel packet retains
its native weights and becomes reciprocal at the period boundary.
Actual finite-field Gauss inputs identify the six-channel packet
with the diagonal Jacobi polynomial of $z^7=x(1-x)$, and its compensated
Abel trace with the affine point-count discrepancy.
The complete cyclic correlation problem also has a uniform constant-eight
stability theorem, proved by an exact signed polynomial certificate.
The arithmetic proofs are maintained in Lean; classical source inputs,
projective geometry and global signal questions are stated separately.
No new zero-free theorem is claimed. The installed local theorem remains
the project's compiled $7/8$ result. Later external theorem statements are
recorded with their versions and have not been imported or proof-audited
here. The former $H_1$ research goal remains paused.
\end{abstract}
\tableofcontents

\section{What is being classified}

A usable case is a tuple
\[
 (F,S,G,n,\mathcal D,\Theta,\mathcal F,\eta,\mathcal T).
\]
Here $F$ is a number field, $S$ the bad places, $G$ the reductive group,
$n$ the covering degree, $\mathcal D$ the covering and splitting data,
$\Theta$ the automorphic representation, $\mathcal F$ a Fourier functional,
$\eta$ the target Hecke character, and $\mathcal T$ the test vectors and
profiles. The cover degree, group rank, residue-character order and
degree of a polynomial input are different parameters.

For a Brylinski--Deligne cover, one retains its quadratic and extension
data, not just $n$. On the cocharacter lattice $Y$ the useful invariants
include
\[
 B_Q(y,z)=Q(y+z)-Q(y)-Q(z),\qquad
 Y_{Q,n}=\{y:B_Q(y,Y)\subset n\ZZ\},\qquad
 n_\alpha=\frac n{\gcd(n,Q(\alpha^\vee))}.
\]
Different root lengths can therefore have different effective orders in
the same cover. Covering-group $L$-groups provide a structural language,
but are not a complete classification of global automorphic forms
or their analytic estimates \cite{Weissman}.

There are four substantially different sources of automorphic input:
\begin{enumerate}
\item The Weil representation and classical quadratic theta series.
Their oscillator and theta-correspondence structure is special to the
double-cover setting.
\item Residual or exceptional theta representations, obtained from
Eisenstein residues. These are closest to the original QRH input.
\item Full metaplectic Eisenstein series and their Fourier coefficients.
They retain spectral parameters and scattering matrices that a residue
may hide.
\item Cuspidal representations and descent constructions. Their summation
formulas can be stronger in some respects, but replacing residual theta by
a cusp form can remove the very residue used to detect a zero.
\end{enumerate}
Products, restrictions and theta lifts create further constructions from
these sources; they must still specify which of these analytic objects
is actually being estimated.

The Fourier functional is another independent choice. Full Whittaker,
semi-Whittaker, degenerate Whittaker, Fourier--Jacobi and descent
coefficients have different supports. A zero full Whittaker model does
not rule out useful degenerate coefficients. Conversely, a nonzero
degenerate coefficient does not establish the scalar Euler factor
needed by a particular detector.

This map covers the principal structural mechanisms, not every
automorphic representation. General root systems of types $D,E,F$,
anisotropic inner forms and ramified representations cannot be honestly
collapsed into a short universal coefficient table.

\section{The cubic baseline and its actual advantages}

The original construction over $\QQ(\sqrt{-3})$ has a cubic theta
coefficient, a sextic character family and cube completion. Its local
coefficient support is $0,1\pmod3$, with the remaining class zero.
The completed ideal index is $c b^3$, with $c$ squarefree; $c,b$ need
not be coprime. The moving character acts on the whole completed index.
After Poisson summation, rows take the form $u a^6$ with $u$ sixth-power-free.
These facts belong together \cite{QRH}.

The mechanism uses the following compatible ingredients.
\begin{enumerate}
\item Explicit theta coefficients whose Gauss factors cancel against
the chosen companion Gauss factors, producing a signed inverse coefficient.
\item A reflected identity for the completed sum, with every local
factor and cusp phase retained.
\item Joint inverse and plain-character witnesses at the same row and
height; the full Hecke family is closed under the twists that occur.
\item Arithmetic moment estimates on that actual family, including
shared-prime zeros and exceptional rows.
\item Prime compensation that cancels the unwanted principal terms
without cancelling the target reciprocal pole.
\end{enumerate}
The ingredient supplied by a better detector is precision and signal
retention. A boundary improvement additionally requires an independent
arithmetic estimate for the resulting signed expression.

In the fixed original Stage-II geometry,
\[
 C_{\mathrm{II}}(s)=s-\frac{11}{16},\qquad
 l_x=\frac{17}{48},\quad b=\frac18,\qquad
 E_{\rm low}=\frac{l_x}{2}+\frac b{12}=\frac3{16}.
\]
The principal exponent at $7/8$ is $3/16$; at $3/4$ it is $1/16$.
Thus the fixed consumer would need a $1/8$ reduction in the error
exponent, with all other errors controlled. Removing only $b/12=1/96$
would leave the boundary $83/96$. This calculation limits that fixed
geometry, not every metaplectic construction.

The full compensation is an alternating sum over every subset of the
chosen prime tuple. It cannot be replaced by its preferred principal
summand. Likewise every numerator, reciprocal $L$-factor, character
phase and nonunit mask in a reflected row remains part of the target.
A numerator-only character estimate is not an estimate for that quotient.

\section{Literature changes that affect the research map}

The review date matters. Three developments change the available inputs.
The October 2 fixed-order large sieve is already a theorem statement,
not the earlier quartic conjecture. The October number-field QRH papers
already pursue changing the base field while retaining cubic theta.
Very recent quartic work treats a complete cusp vector rather than an
imagined scalar coefficient.

Liu's arXiv version 2 states a boundary
$B_\flat=0.874886647039964\ldots$ for all finite-order Hecke
$L$-functions over arbitrary number fields, using an argument over
$K(\zeta_{12})$ and descent \cite{LiuField}. His November 3 Hexagon
version 2 states the later boundary
\[
 B_{\rm ex}=0.87407284766682790777\ldots,\qquad
 178848B^4-507168B^3+522846B^2-231405B+37099=0
\]
with the root in $(0.87407284,0.87407285)$.
The stated mechanism enlarges amplified moment ranges and permits
complete prime removal with a remaining length excess
\cite{LiuLatest}. These are external source statements, not newly
verified local results. Neither is a change to a quartic or higher theta
cover, and neither reaches $3/4$.

The general-field theta input can be inspected directly in the October 4
source: its threefold cover uses a specified central datum, nonzero global
Whittaker functional and fixed bad-place vectors, with exterior support
$c b^3$. Thus passage to fields containing more roots of unity need not
require inventing a new higher-order theta representation
\cite{Primitive}. This field-extension route already exists and should
not be offered as a new original project.

\section{A structural map by group and Fourier model}

The following are unramified local theta-model statements, with the
normalizations and exceptional-character conditions of the cited source.
They do not supply a global detector or estimates at the bad places.
For Kazhdan--Patterson covers of $\GL_r$, the distinguished cases are
$n=r$, or $n=r+1$ with a special covering datum.
For the indicated simply connected groups Gao gives the following list
\cite[Theorems 1.1, 1.2, 5.3, 6.2, 7.1]{Gao}.
\begin{center}
\begin{tabular}{lll}
\toprule
Group & Possible distinguished degrees & Qualification\\
\midrule
$A_r$ & $r+1,r+2$ & stated classification restricted to $n\le r+2$\\
$C_r$, $r\ge2$ & $2r+1,4r-2,4r,4r+2$ & exceptional character matters\\
$B_r$, $r\ge3$ & $2r+1,2r+2$ & exceptional character matters\\
$G_2$ & $7,12$ & specified standard quadratic datum\\
\bottomrule
\end{tabular}
\end{center}
For $C_r$ at $n=4r$, every exceptional character has the distinguished
property. At rank two this makes $\Sp_4$ with degrees $5,6,8,10$
a concrete shortlist. For $D,E,F$ this table makes no classification claim.

A one-dimensional model makes scalar factorization plausible; it does
not make the resulting scalar coefficient a linear-index Gauss sum.
Increasing the rank changes the torus index, determinant height and
the number of summation coordinates. The cost must be computed after
choosing the coefficient slice.

For higher-rank theta with no full Whittaker coefficient, semi-Whittaker
coefficients can be natural replacements. The existing project screen
already examined the $\GL_r$ orbit partition $(n^a,b)$ when $r=an+b$,
and particular Shalika integrals. Repeating that screen is not a new route.
The question for a new integral is its target $L$-factor, all auxiliary
automorphic input, and a usable moving-twist identity.

Descent gives a second way to change the representation instead of
changing a scalar coefficient formula. Friedberg--Ginzburg develop
quartic/sextic descent constructions, with conjectural identifications
in parts of the general picture \cite{FGDescent}; their symplectic
construction includes an unconditional sixfold $\Sp_4$ case while
broader statements retain hypotheses \cite{FGSymplectic}.
A candidate derived this way must be checked for reduction back to the
same cubic input before it is credited with independent analytic power.

\subsection{An all-type screen in the coprime cover range}

For split simply-connected almost-simple groups with the primitive
quadratic datum, Gao--Qiu
give an infinite-family screen, rather than a finite degree list
\cite[Section 3, Table 6, (3.5)]{GaoQiu}. The local field contains
$\mu_n$; retain the tame, unramified and exceptional-character setup.
Writing $e_1,\ldots,e_r$ for the Weyl exponents,
\[
 \dim\operatorname{Wh}_{\psi}(\Theta^{(n)})
   =\frac{\prod_i(n-e_i)}{|W|}
\]
in the following admissible range:
\begin{center}
\begin{tabular}{lll}
\toprule
Type & Admissible $n$ & Dimension-one degree $h+1$\\
\midrule
$A_r$ & $\gcd(n,r+1)=1$ & $r+2$\\
$B_r,C_r$ & $n$ odd & $2r+1$\\
$D_r$ & $n$ odd & $2r-1$\\
$G_2$ & $\gcd(n,6)=1$ & $7$\\
$F_4,E_6$ & $\gcd(n,6)=1$ & $13$\\
$E_7$ & $\gcd(n,6)=1$ & $19$\\
$E_8$ & $\gcd(n,30)=1$ & $31$\\
\bottomrule
\end{tabular}
\end{center}
The final column follows from the formula: each smaller admissible
positive integer is an exponent, whereas exponent duality gives
$\prod_i(h+1-e_i)=|W|$. Above $h+1$ every factor increases.
Thus this column is exhaustive within the indicated range.
This is a consequence of cited literature, not a new Lean theorem
or a global Fourier-coefficient identification.

Outside that range, the earlier even-degree and mixed-root cases
remain separate. Current wavefront-set results also give a systematic
way to select degenerate Fourier models: positive real Satake
parameters give an upper bound with equality in specified cases
\cite{GaoWangWavefront}; for Kazhdan--Patterson and Savin general-linear
covers, highest derivatives determine the wavefront set
\cite{GaoWangZouWavefront}. General covering duality retains conjectural
cases \cite{CoveringBV}. These support a structural screen, not an
analytic estimate for a moving prime-sampled row.

\section{The degree-by-degree landscape}

The field in the next table is the minimal cyclotomic field containing
$\mu_n$. It is not a blanket field-of-definition assertion for every
global construction. Formulas requiring $\mu_{2n}$ can increase the
degree, particularly for even $n$.
\begin{center}
\begin{tabular}{rclp{8.0cm}}
\toprule
$n$ & $[\QQ(\zeta_n):\QQ]$ & source type & Main distinction\\
\midrule
2 & 1 & Weil/theta & Square support and quadratic transformations;
not a direct substitute for a linear Gauss-coefficient detector.\\
3 & 2 & $\GL_2$ theta & Explicit cubic input with sextic twists;
the installed QRH baseline.\\
4 & 2 & $\GL_2$ vector / special $\GL_3$ theta &
Gaussian quartic series have a full cusp-vector reflection;
the scalar higher-rank alternative is a different representation.\\
5 & 4 & $\GL_2$, $\GL_4$, $\Sp_4$ &
The project's $\GL_4$ local model exists; a full global
transformation is unresolved. $\Sp_4$ offers a different rank-two model.\\
6 & 2 & genuine sextic theta / $\Sp_4$ &
Sextic theta is not the cubic theta used with sextic twists.
Descent and coefficient-orbit information are relevant.\\
7 & 6 & special $\GL_6$ / $G_2$ &
$G_2$ gives a rank-two distinguished alternative to increasing
the general-linear rank.\\
8 & 4 & special $\GL_7$ / $\Sp_4$ &
$\Sp_4$ has two root lengths; powers of an order-eight character
descend to quadratic and principal characters.\\
9 & 6 & $\GL_2$ vector / special $\GL_8$ &
Odd-order doubling remains cyclic; large degree and additional
coefficient data must be paid for.\\
10 & 4 & $\GL_2$ vector / $\Sp_4$ &
A rank-two distinguished option with nontrivial fifth-order
characters after squaring.\\
12 & 4 & $\GL_2$ vector / $G_2$ &
Different $G_2$ roots have effective orders $12$ and $4$;
a mixed-order rather than a scalar quartic problem.\\
\bottomrule
\end{tabular}
\end{center}

For $\GL_2$ higher theta, Hecke relations alone need not determine all
coefficient orbits. The quintic pattern has the known orbit
$\{0,3\}$, an undetermined pair $\{1,2\}$ and zero class $4$;
the sextic pattern also has undetermined sectors
\cite{BrokerHoffstein}. These are statements about those constructions,
not a prohibition on all representations of that covering degree.

For quartic theta, Suzuki's square-index evaluations provide information
even when general coefficients are not explicit. David--Dunn--Hamieh--Lin
use them with a quadratic large sieve in their quartic prime-sum method;
their Type-I treatment retains growing square indices and their
conductor costs \cite{DDHL}. This is a viable source of averaged
information, not permission to assign guessed values to other indices.

A fixed integer degree is not a sufficient performance score. One should
compare the root orders, number of genuine torus coordinates, field
degree, unit quotient, bad-prime vector, coefficient height and
transformed conductor together. A degree-twelve rank-two construction
can be structurally simpler than a degree-twelve construction on
$\GL_{11}$, without necessarily having a better analytic bound.

\section{Available character estimates and their exact meaning}

For a fixed order $r\ge3$, a fixed field containing $\mu_r$ and the
fixed family data, de Faveri states
\[
 \Theta_r(A,B)\ll (AB)^\eps
 \left(A+B+A^{1-1/r}B^{2/r}+A^{2/r}B^{1-1/r}\right).
 \tag{6.1}
\]
Both ideal variables are $r$th-power-free. The constants depend on the
fixed field, order and family construction \cite[Theorem 1.1]{deFaveri}.
There is no asserted uniformity in an arbitrarily growing field,
cover degree or additional twist.

For $A=Z^a,B=Z^b$, the exponent supplied by (6.1) is, by elementary
comparison,
\[
 e_r(a,b)=\max(a,b)+
 \frac{\max(0,2\min(a,b)-\max(a,b))}{r}.
 \tag{6.2}
\]
Indeed, if $a\ge b$, the larger mixed exponent is
$a+(2b-a)/r$; the remaining cases follow by symmetry.
At equal lengths $A=B=T$, this is $T^{1+1/r}$:
\[
\begin{array}{c|rrrrrrrr}
r&3&4&5&6&7&8&10&12\\ \hline
1+1/r&4/3&5/4&6/5&7/6&8/7&9/8&11/10&13/12 .
\end{array}
\]
These are exponents of a character-kernel norm, not QRH boundaries.
When $a\ge2b$ the displayed exponent is already $a$ and this particular
advantage of increasing $r$ disappears.

The same paper's root-number coupling is methodologically important:
repeated prime powers are retained as variables and coupled to the
opposite squarefree kernel before estimating. Their conductor can be
much smaller than their index norm. Applying its final scalar bound to
an old squarefree envelope is therefore different from rebuilding an
actual higher-order moment with that coupling retained. Such a rebuild
still requires coefficients independent of the outer moving row.

The earlier Rankin--Selberg analysis proves non-orthogonality in cubic
and quartic families and identifies Gauss bias \cite{FDDH}.
A lower bound exhibiting a biased vector is not a theorem that the
entire excess operator has finite rank. Deleting a few bias vectors
does not establish a sharp bound on the remaining operator.

Cubic cusp-form methods provide another analytic source, including
power-saving prime sums for a fixed form \cite{DunnCusp}. Constants,
levels and the cuspidal hypothesis must survive transfer. Removing the
residual theta component may remove the target zero signal.

\section{What earlier project work already established or attempted}

This review inspected the maintained quintic modules, the prior
metaplectic screen, the October 7 field comparison and the October 8
Gaussian bias-removal note, and the October 9 quintic source bridge.
It searched public and private QRH records for the additional group
and isotropic terminology. This is not an exhaustive audit of every
sibling repository or every equivalent unnamed formulation.

\begin{center}
\begin{tabular}{p{4.1cm}p{9.7cm}}
\toprule
Existing direction & Status relevant to this survey\\
\midrule
Optimal fixed-order sieve & Already recorded locally on October 7;
not a new input discovered by this review.\\
Simple sextic sieve substitution & Earlier complete optimization leaves
the decisive range unchanged. It must not be restarted as a new strategy.\\
Gaussian bias projection & Known norm-induced biased directions were
controlled; the residual image and cross terms were not.
No full residual-operator estimate was proved.\\
Quintic $\GL_2$ scalar adapter & Its known cube-index orbit conflicts
with simultaneous quadratic-reflection and terminal-Gauss conditions.
This excludes that adapter, not quintic theta in general.\\
Quintic $\GL_4$ & Compiled finite local support, exponent and cocycle
algebra; source-dependent linear-index coefficient bridge.
No compiled global theta identification or moving-twist reflection.\\
Completion and inverse formulas & Existing defined arithmetic completion
has no automatic power loss. It does not identify the whole theta slice.\\
Higher-rank/Shalika screen & Already studied selected general-linear
cases and their global auxiliary inputs. A new group requires a new
local and global calculation.\\
\bottomrule
\end{tabular}
\end{center}

In particular, coefficient comparisons and multiplicity recovery do not
supply an estimate for a new metaplectic moment. Their role is to transport
or read out a signal after an actual analytic inequality is available.

\section{A checked cancellation tradeoff in the quintic model}

This is an exact calculation in the existing arithmetic model, with no
global theta identification assumed. The quantified finite domain is
all $v=(a,b,c)\in(\ZZ/5\ZZ)^3$, using its maintained 24-state support.
It is not a scan over finitely many primes.

The symmetric prime-interaction form is
\[
 B(u,v)=2u_1v_1+3(u_1v_2+u_2v_1)
 +4(u_1v_3+u_3v_1)+6u_2v_2
 +8(u_2v_3+u_3v_2)+12u_3v_3.
\]
Among supported residue states, $B(v,v)=0\pmod5$ holds exactly at
\[
 (0,0,0),\quad(1,1,1),\quad(2,2,2),\quad(3,3,3).
\]
This finite classification was checked against every state.
For $v=(1,1,1)$,
\[
 B(mv,nv)=50mn.
\]
The unsymmetrized torus cocycle exponent on this ray is $-25mn$,
also divisible by five. Thus these mixed torus phases cancel for
arbitrary nonnegative integers $m,n$, not just tested prime pairs.

The model's determinant height is $a+2b+3c$. The three nonzero states
above have heights $6,12,18$, squared coefficient powers $1,3,4$,
and defects $5,9,14$. Their first residue-period terms therefore have
substantial height costs. Periodic translates and target-character
powers still have to be retained. The ray is not invariant under all
simple reflections: for example the maintained first reflection sends
$(1,1,1)$ to $(2,3,1)$ modulo five.

Consequently this calculation identifies where a phase-free restriction
is expensive; it does not provide a global reflected packet or an
analytic saving. It has not been promoted as a new Lean theorem.
Its reproducible scope and exact output are in the companion JSON.

\section{Two rank-two families worth distinguishing explicitly}

\subsection{The complete degree-eight $\Sp_4$ coordinate model}

We now compute this case, including its extra central sign.
Work at a tame unramified place whose field contains $\mu_{16}$;
write $q$ for the residue cardinality and $t=q^{-1/2}$.
Number the short coroot first, reversing Gao's numbering in type $C_2$.
Then
\[
 Q(a,b)=a^2-2ab+2b^2,\quad
 Y_{Q,8}=4\ZZ\times4\ZZ,\quad
 Y_{Q,8}^{\mathrm{sc}}=8\ZZ\times4\ZZ.
\]
The exceptional-character values on the two simple root periods are
$q^{-1}$. Thus the value at $(4,0)$ is $\sigma t$, where $\sigma=\pm1$,
and the value at $(0,4)$ is $t^2$.
Both central extensions are distinguished in Gao's classification
\cite[Section 5.2.4]{Gao}. The coordinates below are functional
coordinates in his Section 3.1 convention.

Let $G_k$ be the normalized Gauss sum of $\chi^k$ for an order-eight
character $\chi$, and set
\[
 g=G_1,\quad h=G_2,\quad j=G_3,\quad e=G_4.
\]
All have norm one, $G_{8-k}=G_k^{-1}$ and $e^2=1$.
Normalize $c(0,0)=1$. The complete nonzero rectangle is
\[
\begin{array}{c|c|c}
(a,b)&c(a,b)&Q(a,b)\bmod8\\ \hline
(0,0)&1&0\\
(1,0)&tg&1\\
(0,1)&th&2\\
(1,2)&t^2ge&5\\
(3,1)&t^2hj&5\\
(0,2)&\sigma t^2gje&0\\
(3,3)&t^3hje&1\\
(0,3)&\sigma t^3ghje&2
\end{array}
\]
The other eight rectangle coordinates vanish.

Here is a direct verification of the source recurrence, rather than
a guess based on orbit size. The shifted reflections are
\[
 s_1[a,b]=(1-a+2b,b),\qquad s_2[a,b]=(a,1+a-b).
\]
Put $P_1=2a-2b$, $P_2=-a+2b$, and $n_1=8,n_2=4$.
For an unreduced reflected coordinate $(a',b')$, let
$k=\lfloor a'/4\rfloor$, $\ell=\lfloor b'/4\rfloor$.
Its period multiplier is $\sigma^k t^{k+2\ell}$.
With the trivial Hilbert phase supplied by $\mu_{16}$, Gao's
Corollary 3.7 becomes
\[
 \sigma^k t^{k+2\ell}c(a'\bmod4,b'\bmod4)
 =t^{1-2\lceil P_i/n_i\rceil}
   G_{Q(\alpha_i^\vee)(P_i-1)}^{-1}c(a,b).
\]
Substitution of the displayed table verifies every supported edge;
the remaining edges preserve zero support. In particular the central
$\sigma^k$ cannot be omitted. The kernel checks all 32 edges as exact
degree and Gauss/central-exponent identities.
The source representation theorem is a separately cited input.

The complete signed numerator factors as
\[
\begin{split}
 N_\sigma(X,Y)={}&
 (1+thY)(1+\sigma t^2gjeY^2)\\
 &+tgX(1+teY^2)+t^2hjX^3Y(1+teY^2).
\end{split}
\]
Both unbounded directions are retained in
\[
 \mathcal N_\sigma(X,Y)=
 \frac{N_\sigma(X,Y)}
 {(1-\sigma tX^4)(1-t^2Y^4)}.
\]
The geometric expansion converges absolutely when both period ratios
have modulus less than one. At $X=Y=1$ it converges for $0\le t<1$;
this full two-period convergence is also formalized.
Every natural coordinate pair has the unique form
$(a+4u,b+4v)$ with $0\le a,b<4$ and $u,v\ge0$.
Consequently its squared-coordinate mass, with $x=t^2$, is
\[
 \frac{1+2x+3x^2+2x^3}{(1-x)(1-x^2)}
 =\frac{1+x+2x^2}{(1-x)^2},\qquad 0\le x<1.
\]
This mass differs from a coherent squared sum.

\paragraph{All four interference channels.}
Set $A=gje$, $B=hje/g$.
Grouping by the actual quadratic grade gives, respectively,
\[
 1+\sigma t^2A,\quad
 tg(1+t^2B),\quad
 th(1+\sigma t^2A),\quad
 t^2ge(1+B)
\]
in grades $0,1,2,5$.
Let $S_\sigma(u)$ be the Fourier sum of these four channels against
the standard additive characters of $\ZZ/8\ZZ$.
Finite character orthogonality gives the exact identity
\[
\begin{split}
 E_\sigma&=\frac18\sum_{u\bmod8}|S_\sigma(u)|^2\\
 &=(1+t^2)|1+\sigma t^2A|^2
       +t^2|1+t^2B|^2+t^4|1+B|^2\\
 &=D_8(t)+2\sigma t^2(1+t^2)\Re A+4t^4\Re B,\\
 D_8(t)&=1+2t^2+3t^4+2t^6.
\end{split}
\]
No same-grade cross term is erased.
As a consistency check, removing all cross terms reproduces the
eight-coordinate mass numerator $D_8$.

The classical duplication formula gives
$\chi(4)g/j=he$, hence $B=\chi(4)$.
Here $\chi(4)=\pm1$: the field contains $\mu_8$, so
$\sqrt2=\zeta_8+\zeta_8^{-1}$ belongs to it and the quadratic character
of $2$ is one.
The Gauss--Jacobi identity gives
\[
 A=\frac{J(\chi,\chi^3)}{\sqrt q}.
\]
The automorphism $\zeta_8\mapsto\zeta_8^3$ swaps the two Jacobi
arguments, so fixes this Jacobi sum. Its fixed field is
$\QQ(\sqrt{-2})$, and integrality and the Gauss norm yield
\[
 J(\chi,\chi^3)=a+b\sqrt{-2},\qquad a,b\in\ZZ,\quad a^2+2b^2=q.
\]
These are classical finite-field interpretations of the phases.
The maintained Lean energy theorem instead accepts arbitrary unit
complex Gauss phases with its stated algebraic hypotheses.

\paragraph{Unique compensation with both infinite periods.}
Put
\[
 F_\sigma(u)=\frac{S_\sigma(u)}
 {(1-\sigma t)(1-t^2)},\qquad
 \mathcal E_\sigma=\frac18\sum_{u\bmod8}|F_\sigma(u)|^2
 =\frac{E_\sigma}{(1-\sigma t)^2(1-t^2)^2}.
\]
Averaging the rectangles cancels $\Re A$, but equal averaging of the
complete expressions gives
\[
 \frac{\mathcal E_++\mathcal E_-}{2}
 =\frac{(1+t^2)(D_8+4t^4\Re B)
          +4t^3(1+t^2)\Re A}{(1-t^2)^4}.
\]
The periods reintroduce the phase. There is, however, an exact remedy.

For every $0<t<1$ the unique normalized mixture
$w\mathcal E_++(1-w)\mathcal E_-$ that is independent of the free
real phase $\Re A$ is
\[
 w_+(t)=\frac{(1-t)^2}{2(1+t^2)},\qquad
 w_-(t)=\frac{(1+t)^2}{2(1+t^2)}.
\]
Both weights are strictly positive, sum to one, and satisfy
\[
 \boxed{\quad
 w_+\mathcal E_++w_-\mathcal E_-
 =\frac{D_8(t)+4t^4\Re B}
 {(1+t^2)(1-t^2)^2}.\quad}
\]
To prove uniqueness, the coefficient of $\Re A$ in a general mixture
is
\[
 \frac{2t^2(1+t^2)}
 {(1-t^2)^4}
 \left(2w(1+t^2)-(1-t)^2\right).
\]
Its prefactor is nonzero for $0<t<1$, forcing the displayed $w$.
Conversely substitution cancels that coefficient exactly.
At $t=0$ every mixture is phase independent and uniqueness is not claimed.
The unit rectangle coordinate remains one.
These identities, positivity, uniqueness, and the full coherent
series are maintained Lean theorems.

This gives a precise local design constraint for a higher-cover
construction. It also records a cost: the compensation removes
\emph{all} information carried only by $A$. Unit-coordinate survival
does not establish survival of a target $L$-function zero.
Furthermore one cannot prescribe the signs $\sigma_{\mathfrak p}$
independently at every prime and call the result an admissible global
exceptional-character average. The global character construction,
the intended spherical Fourier functional, reciprocal target factor,
moving-twist conductor and independent signed estimate remain open.
This local theorem changes no certified QRH exponent.

\subsection{The degree-seven and degree-twelve $G_2$ options}

These are structurally distinct from the general-linear screen.
For short and long simple coroots, the quadratic values are $1,3$.
The root periods are therefore $(7,7)$ or $(12,4)$.
The respective minimal cyclotomic field degrees are six and four;
global constructions may require additional roots of unity.

Here is an independent finite check of the unramified orbit count.
Write a cocharacter as $(a,b)$, with $\rho=(5,3)$.
On
\[
 Y/Y_{Q,n}=(\ZZ/n\ZZ)\times
            (\ZZ/(n/\gcd(n,3))\ZZ)
\]
the two shifted reflections are
\[
 s_1[a,b]=(1-a+3b,b),\qquad
 s_2[a,b]=(a,1+a-b).
\]
The full Weyl group has order twelve. Enumerating every coset gives
\[
\begin{array}{c|c|l|c}
 n& |Y/Y_{Q,n}|&\text{orbit sizes}&\text{free orbits}\\ \hline
 7&49&12,6,6,6,6,6,6,1&1\\
 12&48&12,6,6,2,6,6,6,3,1&1 .
\end{array}
\]
This checks these two finite cases; the cited structural theorem,
not extrapolation from this calculation, supplies the all-degree statement.

The degree-twelve case is particularly different: it couples quartic
and twelfth-order root data in two torus variables. A possible analytic
advantage would have to come from this mixed root geometry and its
two-variable functional equations. There is no justification yet for
assigning it the cubic principal exponent or a $3/4$ boundary.
The first test is the same explicit reciprocal-factor test as above,
now on a two-root generating series, with all Weyl terms retained.

Neither option is asserted novel as a representation; both are existing
mathematical objects. Adapting them to a useful zero comparison is a
research possibility, not an established new theorem or a completed
portfolio novelty check.

\section{A phase-free sevenfold $G_2$ model}
\label{sec:G2-seven}

The sevenfold case has a different interference structure from both
models above. Assume the source field contains $\mu_7$ and its residue
characteristic is not seven, as in Gao's tame hypothesis.
Then $q\equiv1\pmod7$ and $t=q^{-1/2}$.
Residue characteristics two and three are not excluded merely by the
cover degree. An order-seven character satisfies $\chi(-1)=1$, so
the normalized Gauss factors satisfy $G_{7-k}=G_k^{-1}$.
Write $g=G_1$, $h=G_2$, $j=G_3$.

Here $Y_{Q,7}=7\ZZ^2$, both root orders are seven, and the central
exceptional character has no extra sign choice.
The complete supported table is
\[
\begin{array}{c|c|c}
(a,b)&c(a,b)&Q(a,b)\bmod7\\ \hline
(0,0)&1&0\\
(1,0)&tg&1\\
(0,1)&tj&3\\
(1,2)&t^2&0\\
(4,1)&t^2&0\\
(6,2)&t^3/h&5\\
(4,4)&t^3h&2\\
(6,5)&t^4&0\\
(2,4)&t^2&0\\
(3,5)&t^3/j&4\\
(2,6)&t^3/g&6\\
(3,6)&t^4&0 .
\end{array}
\]
All other 37 rectangle coordinates vanish.
The same shifted reflections as above, with both periods seven,
verify the root recurrence on every supported edge.
For $(a',b')=s_i[a,b]$, its period multiplier is
$t^{2(\lfloor a'/7\rfloor+\lfloor b'/7\rfloor)}$.
The recurrence multiplier before this reduction is
$t^{1-2\lceil P_i/7\rceil}
G_{Q(\alpha_i^\vee)(P_i-1)}^{-1}$, where
$P_1=2a-3b$ and $P_2=-a+2b$.
Direct substitution checks the full degree and phase table.
These finite checks are kernel theorems. The interpretation as the
source functional is the separate representation-theoretic input
\cite[Corollary 3.7 and Theorem 7.1]{Gao}.

The complete signed polynomial is
\[
\begin{split}
N_7(X,Y)={}&
 (1+t^2X^2Y^4)(1+t^2XY^2+t^2X^4Y)\\
 &+t(gX+jY)\\
 &+t^3(hX^4Y^4+h^{-1}X^6Y^2
             +j^{-1}X^3Y^5+g^{-1}X^2Y^6).
\end{split}
\]
In particular, every nonconstant Gauss phase is a single factor or
its inverse; no mixed Gauss product is being discarded.

The quadratic grades have multiplicities six at zero and one at
each of $1,\ldots,6$. Every pair of supported coordinates in the
same grade has \emph{identical} phase.
Their seven coherent channel coefficients are
\[
\begin{array}{c|ccccccc}
r&0&1&2&3&4&5&6\\ \hline
v_r&1+3t^2+2t^4&tg&t^3h&tj&t^3/j&t^3/h&t^3/g .
\end{array}
\]
For $S(u)=\sum_{r\bmod7}v_r e(ur/7)$, orthogonality therefore gives
\[
\begin{split}
 \frac17\sum_{u\bmod7}|S(u)|^2
 &=(1+3t^2+2t^4)^2+2t^2+4t^6\\
 &=1+8t^2+13t^4+16t^6+4t^8=:E_7(t).
\end{split}
\]
This exact identity holds for \emph{every} unit complex choice of
$g,h,j$. It retains the thirty off-diagonal ordered overlaps in the
grade-zero fiber. There is no unproved Gauss-product relation in it.

Both coherent period directions have ratio $t^2$ at $X=Y=1$.
Their absolutely convergent sum for $0\le t<1$ gives
\[
 F(u)=\frac{S(u)}{(1-t^2)^2},\qquad
 \frac17\sum_{u\bmod7}|F(u)|^2=\frac{E_7(t)}{(1-t^2)^4}.
\]
The channel identity, phase-free energy and full infinite-period
convergence are maintained Lean theorems.

The squared-coordinate mass has a distinct closed form. For $x=t^2$,
the finite numerator is
\[
 1+2x+3x^2+4x^3+2x^4=(1+x)^2(1+2x^2).
\]
Each coordinate period contributes squared ratio $x^2$, so the mass
over all natural coordinate pairs is
\[
 M_7(x)=\frac{1+2x^2}{(1-x)^2},\qquad0\le x<1.
\]
The unique residue/period indexing and this convergent mass are also
formalized. As algebraic consequences of the three established mass
formulas, at the same formal $0<x<1$,
\[
 M_{12}(x)-M_7(x)=\frac{x}{1-x}>0,\qquad
 M_{8}(x)-M_{12}(x)=\frac{x^2}{(1-x)^2}>0.
\]
These comparisons concern the normalized coordinate models; their
minimal global fields and conductors differ.

This calculation supplies a useful classification criterion for the
next analytic work: count the \emph{phases inside each grade}, rather
than just the number of coordinates or the cover degree.
The sevenfold total energy retains no information about $g,h,j$.
The cubic observable below recovers the phase $G_1G_2G_4$ that this
second moment misses. The constructive grade-zero overlaps also show why
a smaller coordinate mass is not itself a smaller coherent main term.

\section{A principal-preserving cubic readout in the sevenfold model}
\label{sec:G2-seven-cubic}

Put $e_7(r)=\exp(2\pi i r/7)$ and let brackets denote the complete
average over $a\in\ZZ/7\ZZ$. Retain the entire principal polynomial
\[
 P_0=(1+t^2X^2Y^4)(1+t^2XY^2+t^2X^4Y).
\]
Project the complete source-coordinate sum onto grades $\{0,1,2,4\}$
and $\{0,3,5,6\}$ respectively. The resulting carriers are exactly
\[
\begin{split}
 S_R(a)&=P_0+tgX e_7(a)+t^3hX^4Y^4e_7(2a)
                     +t^3j^{-1}X^3Y^5e_7(4a),\\
 S_N(a)&=P_0+tjY e_7(3a)+t^3h^{-1}X^6Y^2e_7(5a)
                     +t^3g^{-1}X^2Y^6e_7(6a).
\end{split}
\]
Every grade-zero coordinate remains in each carrier. These formulas
are theorems about the full 49-state arithmetic model, rather than
definitions that omit the other principal coordinates.

\textbf{Cubic readout theorem.}
For arbitrary nonzero complex $g,h,j$ and arbitrary complex $t,X,Y$, set
\[
 \Phi_R(a)=1+\frac{(2P_0-S_R(a))^3-P_0^3}{6},\qquad
 \Phi_N(a)=1+\frac{(2P_0-S_N(a))^3-P_0^3}{6}.
\]
Then, with $A=gh/j$ and $z=t^7X^8Y^9$,
\[
 \langle\Phi_R\rangle=1-Az,\qquad
 \langle\Phi_N\rangle=1-A^{-1}z,
\]
and hence
\[
 \langle\Phi_R\rangle\langle\Phi_N\rangle
       =1-(A+A^{-1})z+z^2.
\]
This is the product of two separately averaged probes; identifying
it with a same-parameter average would require additional cross terms.
If $|g|=|h|=|j|=1$, its trace coefficient is $2\Re A$.

\emph{Proof.}
For either frequency triple $(r_1,r_2,r_3)=(1,2,4)$ or $(3,5,6)$,
expand
\[
 C(a)=u+A_1e_7(r_1a)+A_2e_7(r_2a)+A_3e_7(r_3a).
\]
Orthogonality applies to each of the $4^3=64$ ordered terms in $C(a)^3$.
Among triples drawn from $\{0,r_1,r_2,r_3\}$, the zero-frequency
triples are exactly $(0,0,0)$ and the six permutations of
$(r_1,r_2,r_3)$. In particular, none with exactly one or two
nonprincipal factors survives. Therefore
\[
 \langle C^3\rangle=u^3+6A_1A_2A_3 .
\]
The generic identity holds for arbitrary complex $u,A_1,A_2,A_3$.
Apply it with $u=P_0$ and negate each nonprincipal coefficient:
\[
 \left\langle(2P_0-S_R)^3\right\rangle=P_0^3-6Az,\qquad
 \left\langle(2P_0-S_N)^3\right\rangle=P_0^3-6A^{-1}z.
\]
Both products have degree $t^{1+3+3}$ and torus bidegree
$(1,0)+(4,4)+(3,5)=(0,1)+(6,2)+(2,6)=(8,9)$.
Subtraction of $P_0^3$ and division by six give the two means.
Their product gives the stated polynomial. Under unit inputs,
$A^{-1}=\overline A$, which proves the real-trace assertion.
All finite orthogonality, resonance, source-coordinate, compensation
and paired identities in this proof are maintained Lean theorems.

\textbf{Full coherent periods.}
Write
\[
 D=(1-t^2X^7)(1-t^2Y^7),\qquad F_R(a)=S_R(a)/D,
 \quad F_N(a)=S_N(a)/D.
\]
For $|t^2X^7|<1$ and $|t^2Y^7|<1$, the complete double series is
\[
 F_R(a)=\sum_{m,n\ge0}(t^2X^7)^m(t^2Y^7)^nS_R(a),
\]
and likewise for $N$. Absolute convergence follows from the product
of the two norm-geometric series, so multiplication of their sums
is justified. Both series and their exact sums are formalized.
Whenever $D\ne0$, replacing $S_R,S_N$ in the probes by $DF_R,DF_N$
gives exactly the same readout: both period factors are compensated
\emph{before} cubing.

\textbf{Exact range of this absolute period box.}
For the physical scale $t=q^{-1/2}$, write $|X|=q^u$ and $|Y|=q^v$.
The two strict norm bounds become $u,v<1/7$. To have $|z|=q^{-\sigma}$
requires
\[
 8u+9v=\frac72-\sigma.
\]
Such $u,v$ exist if and only if $\sigma>15/14$. Necessity follows
from $8u+9v<17/7$; sufficiency follows by choosing
$u=v=(7/2-\sigma)/17$. This exact scalar equivalence is formalized.
It limits this absolute series representation, not the algebraic
compensated identity or every possible global continuation.

For actual normalized order-seven Gauss sums, the classical
Gauss--Jacobi formula interprets $A$ as
$J(\chi,\chi^2)/\sqrt q$. That classical arithmetic interpretation is
separate from the kernel-checked arbitrary-phase theorem above.
A compatible global Hecke character, spherical/global Fourier
functional and independent signed estimate have not been constructed
here. In particular, the local reciprocal-factor shape does not
change the installed $7/8$ boundary. Priority for this particular
combined readout has not been established.

\section{The complete cubic spectrum and its two arithmetic sectors}
\label{sec:G2-seven-spectrum}

The residue-only readout selects a quadratic CM sector, but the full
cubic moment carries more arithmetic. Let $v_0=u$ and $v_1,\ldots,v_6$
be arbitrary complex amplitudes, and put
$S(a)=\sum_{r=0}^6v_re_7(ar)$. Complete orthogonality gives the
following exact theorem:
\[
\begin{split}
 \langle S^3\rangle={}&u^3+
 6u(v_1v_6+v_2v_5+v_3v_4)\\
 &+3(v_1^2v_5+v_1v_3^2+v_2^2v_3
              +v_2v_6^2+v_4^2v_6+v_4v_5^2)\\
 &+6(v_1v_2v_4+v_3v_5v_6).
\end{split}
\]
This includes all $7^3$ ordered terms before averaging. Exactly 49
resonate: the principal triple, eighteen terms containing one principal
factor, and thirty fully nonprincipal terms. No principal cross term
has been removed by centering.

For the complete source-coordinate model, take $u=P_0$ and the six
nonprincipal torus amplitudes from Section~\ref{sec:G2-seven-cubic}.
Define
\[
 A=gh/j,\quad B=g^2/h,\quad C=gj^2,\quad D=h^2j.
\]
The eight sorted nonprincipal triple types are the entire spectrum:
\[
\begin{array}{c|c|c|c}
(r,s,k)&\text{multiplicity}&\text{phase}&\text{torus weight}\\ \hline
(1,1,5)&3&B&t^5X^8Y^2\\
(1,2,4)&6&A&t^7X^8Y^9\\
(1,3,3)&3&C&t^3XY^2\\
(2,2,3)&3&D&t^7X^8Y^9\\
(2,6,6)&3&B^{-1}&t^9X^8Y^{16}\\
(3,5,6)&6&A^{-1}&t^7X^8Y^9\\
(4,4,6)&3&C^{-1}&t^9X^8Y^{16}\\
(4,5,5)&3&D^{-1}&t^9X^{15}Y^9 .
\end{array}
\]
If $\mathcal C$ is the sum of these eight weighted entries, including
their multiplicities, the complete physical identity is
\[
 \left\langle S_{\rm full}^3\right\rangle
 =P_0^3+6P_0(2t^4X^3Y^6+t^6X^{10}Y^6)+\mathcal C.
\]
Both the arbitrary-amplitude identity and this full 49-state
specialization are kernel-checked theorems.

\emph{Proof and completeness.}
Expand the cube and sum each term over $a$. A term survives precisely
when $r+s+k=0$ modulo seven. With a zero index this gives the principal
triple or an opposite pair and its six ordered placements. Without a
zero index, sorting gives exactly the eight triples in the table.
Their permutation counts are three or six as displayed. Substitution
of the actual coordinate amplitudes gives the stated torus exponents
and the principal cross term. The finite classification of every
sorted nonzero triple is itself formalized, rather than inferred
from a selection of examples.

There is a second exact classification. Scale each triple by
$m\in(\ZZ/7\ZZ)^\times$ and compare multiplicities of its entries.
The triples $(1,2,4)$ and $(3,5,6)$ form one orbit, with stabilizer
$\{1,2,4\}$. The other six triples form a disjoint orbit with trivial
stabilizer. For a triple $T$, set
\[
 h_T(m)=\sum_{r\in T}\left\{\frac{mr}{7}\right\}-1.
\]
The sum is exactly integral; no floor or rounding error is hidden.
Its values are zero or one, and $h_T(m)+h_T(-m)=1$.
For $m=1,\ldots,6$, representatives have vectors
\[
 h_{124}=(0,0,1,0,1,1),\qquad
 h_{115}=(0,0,0,1,1,1).
\]
The stabilizers, orbit coverage and disjointness, exact integrality,
conjugation and both representative vectors are all kernel-checked.

The arithmetic interpretation uses a separate classical input.
For actual Gauss sums, every nondegenerate level-seven Jacobi ratio
is a product $G_rG_sG_{-r-s}$, so the table covers all such ratios.
The formal-multiset field and fractional-part rules in
\cite[Sections 2.1--2.3]{Watkins} associate the first orbit with the
quadratic CM field $\QQ(\sqrt{-7})$, and the second with the degree-six
CM field $\QQ(\zeta_7)$, before unitary normalization. The displayed
indices are those after the Tate shift giving weight one.
These are classical interpretations of the proved finite data,
not Lean field-identification theorems. In particular the natural
field of a motive is not a claim that every value at every prime
has that exact field degree.

This gives a concrete reason to retain higher correlations when
categorizing possible extensions: one cover contains more than the
quadratic sector selected by the first probe. The six remaining
terms carry nonconstant infinity-type data. Their unequal physical
torus weights cannot be silently replaced by an unweighted Galois
trace. A compatible global Fourier construction must keep those
weights, the source sign convention, the bad prime seven and the
archimedean factors. Such a construction and an independent signed
analytic estimate remain open here; the installed boundary stays
$7/8$. Priority for this combined source-coordinate spectrum is
not established.

\section{All Jacobi readouts and a coupled boundary packet}
\label{sec:G2-seven-jacobi}

Number the eight preceding triple rows by $i=0,\ldots,7$, with
phases $\phi_i$ and native weights $w_i=t^{d_i}X^{a_i}Y^{b_i}$.
Let $m_i=6$ for rows 1 and 5, and $m_i=3$ otherwise.
The mask for row $i$ keeps grade zero and the distinct nonzero
grades in that row. Its full source projection $S_i(a)$ includes
all six principal coordinates. Define
\[
 Q_i(a)=1+\frac{(2P_0-S_i(a))^3-P_0^3}{m_i}.
\]
For every row and every nonzero $g,h,j$, the exact theorem is
\begin{equation}
 \langle Q_i\rangle=1-\phi_iw_i .
 \label{eq:all-jacobi-readouts}
\end{equation}
All eight source masks and their compensation are kernel checked.

\emph{Proof.} No nonzero opposite pair survives in a mask, and its
only fully nonprincipal zero-sum triple is the selected row, with
multiplicity $m_i$. The complete cubic identity therefore gives,
for arbitrary amplitudes,
\[
 \left\langle
   \left(v_0-\sum_{r\ne0,\ r\ {\rm kept}}v_re_7(ar)\right)^3
 \right\rangle=v_0^3-m_i\prod_{r\in T_i}v_r .
\]
Repeated entries of $T_i$ occur repeatedly in the product.
Taking $v_0=P_0$ and substituting the actual seven torus amplitudes
proves the formula, retaining every native weight.

Write $\Delta=(1-t^2X^7)(1-t^2Y^7)$.
Replace $S_i$ by $\Delta S_i^{\rm coh}$ to define the coherent
probe $Q_i^{\rm coh}$. For $\Delta\ne0$ it equals $Q_i$
pointwise. In the domain $|t^2X^7|<1$, $|t^2Y^7|<1$, the
existing complete two-period summation theorem identifies
$S_i^{\rm coh}=S_i/\Delta$ with the convergent infinite series.
Both factors are compensated before cubing.

The extracted native roots and Euler probes are
\[
 \lambda_i=\langle Q_i\rangle-1=-\phi_iw_i,\qquad
 E_i=2-Q_i,\qquad \langle E_i\rangle=1-\lambda_i.
\]
The sign matters. Gao's local integral uses ordinary Gauss sums.
Their norm/trace-compatible Hasse--Davenport lifting has sign
$(-1)^{n-1}$; modified negative Gauss sums satisfy a power law
\cite[Proposition 2.12]{KubertLichtenbaum}.
Our formal lifting theorem explicitly assumes
$-v'_r=(-v_r)^n$ for every positive grade $r$.
It proves $\lambda_i(v')=\lambda_i(v)^n$ because
$-\prod_{r\in T_i}v'_r=\prod_{r\in T_i}(-v'_r)$, a product
of three factors. Both principal amplitudes remain arbitrary.
Actual finite-field lifting, compatible torus scaling and choices
of characters are separate classical inputs, not new Lean theorems.

Put $W=t^7X^8Y^9$ and $R=t^2X^7$. The native roots satisfy
\begin{align*}
 \lambda_0\lambda_3\lambda_6&=\lambda_1W^2,&
 \lambda_2\lambda_4\lambda_7&=\lambda_5W^2,\\
 \lambda_0\lambda_4&=W^2,&
 R\lambda_2\lambda_6&=W^2,&
 \lambda_3\lambda_7&=RW^2.
\end{align*}
For the first determinant, $BD/C=A$ and the weights multiply to
$W^3$; the second uses $C/(BD)=A^{-1}$. Direct substitution
proves every displayed equality without dividing by $R$ or $W$.
All are kernel checked with the original weights.

At $R=1$, the three pairs have product $W^2$. Define
\[
 {\cal P}(U,V)=
 \prod_{(i,k)\in\{(0,4),(2,6),(3,7)\}}
       (U-\lambda_iV)(U-\lambda_kV).
\]
The exact homogeneous reciprocity theorem is
\begin{equation}
 {\cal P}(W^2V,U)=W^6{\cal P}(U,V).
 \label{eq:native-packet-reciprocity}
\end{equation}
Indeed, for $ab=W^2$,
\[
 (W^2V-aU)(W^2V-bU)=W^2(U-aV)(U-bV).
\]
Multiply this identity for all three pairs. The proof includes
$W=0$ and keeps every product cross term. This packet is a
product of separately averaged probes, not the average of a
same-twist product.

There is a nonvacuous joint Abel limit. Suppose
$t^2X_0^7=t^2Y_0^7=1$, and take $X=rX_0,Y=rY_0$ with
$0\le r<1$. Both period ratios equal $r^7$; both full series
converge absolutely, and $\Delta=(1-r^7)^2\ne0$.
Continuity of the compensated finite polynomial gives
\[
 \lim_{r\to1^-}\langle Q_i^{\rm coh}(rX_0,rY_0)\rangle
 =1-\phi_iw_i(t,X_0,Y_0).
\]
Lean explicitly verifies that this left-neighborhood filter is
nontrivial. Subtract one and multiply the six limiting factors:
the full coherent packet has Abel limit $\mathcal P$, satisfying
\eqref{eq:native-packet-reciprocity}. No divergent period sum is
evaluated at its pole. A separate fixed-$Y$ filtered-extension
theorem is not claimed nonvacuous if its domain is empty.

For $u\ne0$, a canonical common-scale corner is
\[
 t=u^7,\qquad X_0=Y_0=u^{-2}.
\]
Every row satisfies $7d_i=15+2(a_i+b_i)$, so both boundary
equations hold and every $w_i=u^{15}$. Thus
\begin{equation}
 \lim_{r\to1^-}
 \frac{\langle Q_i^{\rm coh}(ru^{-2},ru^{-2})\rangle-1}{u^{15}}
 =-\phi_i .
 \label{eq:normalized-jacobi-abel}
\end{equation}
This all-row normalized limit is kernel checked. It recovers
the degree-six sector as well as the quadratic pair.

The cost is explicit. With $u=q^{-1/14}$, normalization divides
by $q^{-15/14}$. Every individual row also satisfies the sharp
scalar absolute period-box equivalence
\[
 \exists\,\xi,\eta<1/7:\quad
 a_i\xi+b_i\eta=d_i/2-s
 \quad\Longleftrightarrow\quad s>15/14 .
\]
For necessity weight the two strict inequalities by $a_i,b_i>0$.
For sufficiency take
$\xi=\eta=(d_i/2-s)/(a_i+b_i)$ and use the exponent identity.
The equivalence is formalized for all eight rows.

These results expose all Jacobi channels and their coupled
reciprocal boundary packet. The next section identifies the
six-channel packet with an actual finite-field Jacobi polynomial.
An admissible global metaplectic Fourier construction and a
uniform estimate after normalization remain separate questions.
No independent signed character estimate or QRH gain is proved;
the installed boundary stays $7/8$.
Priority for this combined construction is not established.

\section{An actual genus-three curve inside the sevenfold readout}
\label{sec:G2-seven-curve}

The six-channel sector admits a concrete finite-field interpretation.
Let $F$ be any finite field with $q$ elements, $\chi:F\to\mathbb C$
a multiplicative character of exact order seven, and $\psi$ a
primitive additive character. Characters are zero on zero. Choose
$\rho\ne0$ with $\rho^2=q$, and put
\[
 G_r=\frac{\sum_{x\in F}\chi^r(x)\psi(x)}{\rho},\qquad
 J_r=\sum_{x\in F}\chi^r(x)\chi^r(1-x),\qquad 1\le r\le6.
\]
Use the actual inputs $g=G_1,h=G_2,j=G_3$ in the source model.
All statements below concerning these finite sums and their
compensated source consumers are kernel checked.

First, $G_{7-r}=G_r^{-1}$ and all six entries are nonzero.
The exact Gauss--Jacobi convolution gives
\[
 \frac{J_r}{\rho}=\frac{G_r^2}{G_{2r}},
\]
with character exponents reduced modulo seven. Consequently the
six source phases, in character-exponent order, are
\[
\begin{array}{c|rrrrrr}
r&1&2&3&4&5&6\\ \hline
\text{source row }i&0&3&2&6&7&4\\
J_r/\rho&B&D&C&C^{-1}&D^{-1}&B^{-1}.
\end{array}
\]
These are actual finite-field sums, rather than formal Gauss
variables satisfying an unproved identification.

\emph{Proof of the Gauss conversion.}
Expand the product of two Gauss sums and group it by the additive
sum $t$. For $t\ne0$, substitute $x=tu$ to obtain
$\chi(t)\phi(t)J(\chi,\phi)$. At $t=0$ the full fiber is
$\phi(-1)\sum_x(\chi\phi)(x)=0$ when $\chi\phi$ is nonprincipal.
Thus $G_{\rm raw}(\chi)G_{\rm raw}(\phi)
=J(\chi,\phi)G_{\rm raw}(\chi\phi)$, with the zero fiber retained.
For a nontrivial power of $\chi$, the inverse Gauss product is $q$,
since $\chi(-1)=1$. Dividing by $\rho^2=q$ gives the reciprocal
entries. Substitution proves the displayed six-phase table.
All three steps, including the zero fiber, are formalized.

Consider the affine curve
\[
 C_{\rm aff}:\quad z^7=x(1-x),\qquad
 N_{\rm aff}=\#C_{\rm aff}(F).
\]
Its complete character count is
\begin{equation}
 N_{\rm aff}=q+\sum_{r=1}^6J_r.
 \label{eq:curve-affine-jacobi}
\end{equation}
No root-counting hypothesis is assumed in this theorem.

\emph{Proof.} At $a=0$ there is one solution of $z^7=a$.
For $a\ne0$, cyclicity of $F^\times$ shows that $a$ is a seventh
power exactly when $\chi(a)=1$. Each nonempty fiber has seven
elements, by the kernel size of the seventh-power homomorphism.
The complete geometric sum therefore gives
\[
 \#\{z\in F:z^7=a\}=1+\sum_{r=1}^6\chi^r(a)
\]
for every $a$, including zero. Sum this identity at $a=x(1-x)$.
Multiplicativity and interchange of the two finite sums give
\eqref{eq:curve-affine-jacobi}. Lean proves the cyclic-group
argument and the fiber size, not only the final substitution.

Now choose any auxiliary scale $a\ne0$ and set
$t=a^7$, $X_0=Y_0=a^{-2}$. For the native source packet from
Section~\ref{sec:G2-seven-jacobi}, the exact theorem is
\begin{equation}
 {\cal P}_{\rm source}\left(U,\frac{\rho V}{a^{15}}\right)
   =\prod_{r=1}^6(U+J_rV).
 \label{eq:curve-source-packet}
\end{equation}
The right-hand side contains neither $\psi$, nor $\rho$, nor $a$.
Thus the conversion removes those auxiliary choices after the
stated rescaling.

\emph{Proof.} Every native weight is $a^{15}$. Its extracted
root is $-a^{15}\phi_i$. Each rescaled factor is therefore
$U+\rho\phi_iV$. The six-phase table identifies these with
$U+J_rV$, and their complete product proves
\eqref{eq:curve-source-packet}. This is a theorem about the
actual source packet, preserving every factor and cross term.

There is also a kernel-checked source-to-count Abel theorem:
\begin{equation}
 \lim_{s\to1^-}\sum_{r=1}^6
 \frac{\langle Q_{\iota(r)}^{\rm coh}(sa^{-2},sa^{-2})\rangle-1}
      {a^{15}}
 =\frac{q-N_{\rm aff}}{\rho},
 \label{eq:curve-source-trace}
\end{equation}
where $\iota=(0,3,2,6,7,4)$ is the row table above.
Both complete period ratios are $s^7$ in the approach region;
the limiting filter is nontrivial. Apply the six previously
proved normalized Abel limits, sum their negative Jacobi phases,
and use \eqref{eq:curve-affine-jacobi}. The full principal
polynomial and both period compensations remain in each probe.

The smooth-projective interpretation is a separate classical
input. This cyclic Fermat quotient has genus three and one
rational point above infinity. In odd characteristic,
$Y=1-2x$, $X=z$ gives the equivalent affine model
$Y^2=1-4X^7$. In characteristic two the Kummer model instead
reads $x^2+x=z^7$; the counting theorem above still applies.
Characteristic seven is excluded by the tame source datum:
an order-seven character cannot exist on its finite-field unit
group, since seven does not divide $q-1$.
For a rational prime $p\ne7$, the least admissible residue degree
is $f=\operatorname{ord}_7(p)$; it is $1,3,6,3,6,2$ for residues
$1,2,3,4,5,6$, respectively. This does not yet identify the
rational-prime descent of the packet.

The classical Jacobi-sum zeta formula
\cite[Introduction, equation (2)]{AokiQuotients} gives
\[
 Z(C/F,T)=
 \frac{\prod_{r=1}^6(1+J_rT)}{(1-T)(1-qT)} .
\]
Its sign can also be checked from norm-compatible lifting:
$J_{r,n}=(-1)^{n-1}J_r^n$, so
$\#C(F_{q^n})=q^n+1-\sum_r(-J_r)^n$.
This projective geometry and the full lifting/zeta formalism
are cited classical inputs; they are not new Lean declarations
here. The newly formalized bridge is the exact identification
of this finite Jacobi polynomial and affine trace with the full
compensated metaplectic source observation.

This supplies a specific geometric target for further global
work, beyond a polynomial having a suggestive reciprocal shape.
It is not an identification of an admissible global metaplectic
Fourier coefficient. With physical $a=q^{-1/14}$ and
$\rho=\sqrt q$, the trace normalization still costs $q^{15/14}$;
the unnormalized polynomial variable in
\eqref{eq:curve-source-packet} rescales by $q^{11/7}$.
No uniform analytic saving or improvement of $7/8$ follows.
The curve and its zeta formula are classical; priority for this
precise source-observation bridge is not established.

\section{Frobenius compression and the exact vanishing branch}
\label{sec:G2-seven-frobenius}

The preceding identification can now be followed through every
characteristic residue modulo seven. These assertions are proved in
the maintained \texttt{G2SevenCurveFrobenius} module, including their
consumer in the actual compensated source packet.

Let $p$ be the characteristic of $F$, retain the order-seven
character, and write $J_r=J(\chi^r,\chi^r)$. The fundamental identity
is
\[
 J(\chi^p,\phi^p)=J(\chi,\phi),\qquad J_{pr}=J_r,
 \tag{F.1}
\]
where subscripts are reduced modulo seven. For the first equality
no order-seven hypothesis is needed.

\emph{Proof.} Frobenius $x\mapsto x^p$ is a permutation of the
finite source field. Multiplicativity gives
$\chi^p(x)=\chi(x^p)$, and in characteristic $p$,
$(1-x)^p=1-x^p$. Reindex the complete Jacobi sum by Frobenius.
Apply the result to $\chi^r$ and reduce its exponents modulo
$\operatorname{ord}(\chi)=7$. No Frobenius operation is applied to
the complex target.

Let $f=\operatorname{ord}_7(p)$, and let $\mathcal R_p$ contain one
exponent from each multiplication-by-$p$ orbit. Put
\[
 R_p(U,V)=\prod_{r\in\mathcal R_p}(U+J_rV).
\]
The complete classification is
\[
\begin{array}{c|c|l|l}
p\bmod7&f&\text{orbits on }1,\ldots,6&\mathcal R_p\\ \hline
1&1&(1),(2),(3),(4),(5),(6)&1,2,3,4,5,6\\
2&3&(1,2,4),(3,6,5)&1,3\\
3&6&(1,3,2,6,4,5)&1\\
4&3&(1,4,2),(3,5,6)&1,3\\
5&6&(1,5,4,6,2,3)&1\\
6&2&(1,6),(2,5),(3,4)&1,2,3 .
\end{array}
\]
In particular, $\deg R_p=6/f$. The exact identities are
\begin{align}
 \prod_{r=1}^6(U+J_rV)&=R_p(U,V)^f, \tag{F.2}\\
 \sum_{r=1}^6(-J_r)^n&=
 f\sum_{r\in\mathcal R_p}(-J_r)^n\qquad(n\ge0). \tag{F.3}
\end{align}
Both are kernel checked for arbitrary complex values invariant under
the displayed permutation; (F.1) discharges that invariance for the
actual finite-field sums. The proof lists all six residue classes,
follows each complete orbit, and expands the full product or sum.
There is no norm inequality and no discarded cross term.
The power-trace theorem quantifies over every natural $n$.

Substitution in \eqref{eq:curve-source-packet} proves the actual
source theorem
\begin{equation}
 {\cal P}_{\rm source}\left(U,\frac{\rho V}{a^{15}}\right)
       =R_p(U,V)^f .
 \tag{F.4}
\end{equation}
The primitive additive character, arbitrary nonzero $a$ and
$\rho^2=q$ retain their earlier domains. Thus the six-channel
source has six, two, one, two, one or three distinct orbit factors
in characteristic classes $1,2,3,4,5,6$, respectively. These are
exact orbit factors; further accidental equality is not excluded.

There is also an unconditional branch without any character:
\begin{equation}
 q\not\equiv1\pmod7
 \quad\Longrightarrow\quad
 \#\{(x,z)\in F^2:z^7=x(1-x)\}=q .
 \tag{F.5}
\end{equation}
Indeed, $\gcd(7,q-1)=1$, so seventh powering is a bijection on
$F^\times$, and it extends to a bijection on $F$ by fixing zero.
Every $x$ has exactly one corresponding $z$. Lean constructs the
power equivalence on units, proves uniqueness of every fiber
including zero, and sums these fibers. It also derives the
coprimality from the displayed residue condition.
Equation (F.5) includes characteristic seven as an affine counting
identity; it makes no tame-source or smooth-projective assertion
there. In characteristic two the original Kummer model is used,
without dividing by two.

The field in (F.2)--(F.4) already supports an order-seven character.
For $q=p^f$, this is an extension-field factorization, and its $f$th
power must not be mistaken for the rational-prime Euler numerator.
That descent must also use the counts over every $p^n$:
(F.5) treats the nonsplit extension degrees, while compatible
Jacobi lifting and the point above infinity are the remaining
classical geometric inputs. Frobenius compression alone does not
prove that descent in Lean.

This classification is useful for screening generalizations:
cover degree, the number of Gauss coordinates, and the number of
independent orbit factors are different data. Here the exact
source-to-curve bridge retains the principal compensation and
both infinite periods, while Frobenius reduces its effective
orbit count. It supplies no independent global signed estimate.
The costs $q^{15/14}$ and $q^{11/7}$ above are unchanged, as is
the installed $7/8$ boundary. Frobenius and Fermat-quotient
factorizations are classical; priority for this exact compensated
source consumer has not been established.


\section{The complete cubic source reads the Fermat septic}
\label{sec:G2-fermat}

The eight masked readouts have a joint geometric consumer. They
recover all thirty nondegenerate Jacobi contributions to the Fermat
septic, with their original multiplicities, and also assemble to a
single full-source cubic observation. This section states the actual
finite-field and source identities proved in
\texttt{G2SevenFermatCount} and \texttt{G2SevenFermatReadout}.

Retain a finite field $F$ of cardinality $q$, an order-seven complex
character $\chi$, a primitive additive character $\psi$, and
$\rho\ne0$ with $\rho^2=q$. Write
\[
 J_{r,s}=J(\chi^r,\chi^s),\qquad
 S_{30}=\sum_{\substack{1\le r,s\le6\\r+s\not\equiv0\ (7)}}J_{r,s}.
\]
The sum has thirty ordered pairs. Define the actual counts
\[
 N_{\rm aff}=\#\{(x,y)\in F^2:x^7+y^7=1\},\qquad
 N_F=N_{\rm aff}+\#\{z\in F:z^7=-1\}.
\]
The second count is the sum of the two canonical projective normal
forms: $Z=1$, or $(z:1:0)$ at infinity. Every projective point has
exactly one of these representatives. The Lean count is explicitly
this finite chart sum; no projective-geometry library theorem is
assumed in its arithmetic proof.

\paragraph{Complete counting theorem.}
The exact identities are
\begin{align}
 N_{\rm aff}&=q+\sum_{r,s=1}^6J_{r,s}
             =q-6+S_{30}, \label{eq:fermat-affine}\\
 \#\{z:z^7=-1\}&=7,\qquad N_F=q+1+S_{30}.
 \label{eq:fermat-complete}
\end{align}

\emph{Proof.} The previously proved root-count law gives
$R(a)=1+\sum_{r=1}^6\chi^r(a)$, including $a=0$.
The exact seventh-power pushforward yields
\[
 N_{\rm aff}=\sum_{a\in F}R(a)R(1-a).
\]
Both linear character sums vanish by nonprincipality and the
permutation $a\mapsto1-a$. The constant contributes $q$, while
the full product gives all thirty-six $J_{r,s}$.

For $s=7-r$, the correction is $J_{r,7-r}=-1$.
Here is the endpoint-preserving argument formalized in Lean.
For any nonprincipal $\eta$ with $\eta(-1)=1$, and
$\eta\phi=1$, the map
\[
 f(x)=
 \begin{cases}1,&x=1,\\ x/(x-1),&x\ne1\end{cases}
\]
is an involution of the entire field. Direct multiplication,
including $x=1$, gives
\[
 \eta(x)\phi(1-x)=\eta(f(x))-\boldsymbol1_{\{x=1\}}.
\]
Its sum is $-1$. Apply this to $\eta=\chi^r$,
$\phi=\chi^{7-r}$. Odd order gives $\eta(-1)=1$.
Thus exactly six inverse pairs contribute $-6$.
Finally $\chi(-1)=1$, so the same complete fiber law gives seven
solutions at infinity. This proves both displayed identities,
with neither zero fibers nor infinity omitted.

If $q\not\equiv1\bmod7$, seventh powering is bijective instead.
Then $N_{\rm aff}=q$ and $N_F=q+1$, with no character assumption.
This chart-count identity includes characteristic seven; the
smooth degree-seven curve and tame source interpretation use
characteristic different from seven. No division by two occurs.

\paragraph{All thirty terms match the source.}
For every nondegenerate pair, the proved Gauss convolution gives
\[
 \frac{J_{r,s}}{\rho}=\frac{G_rG_s}{G_{r+s}}.
\]
The denominator is nonzero because $\chi^{r+s}$ is nonprincipal.
Reduce exponents modulo seven and use the proved reciprocal
Gauss entries. With the eight phases $\phi_i$ and multiplicities
$m_i$ from Section~\ref{sec:G2-seven-jacobi}, this gives
\begin{equation}
 \frac{S_{30}}{\rho}=\sum_{i=0}^7m_i\phi_i,\qquad
 \frac{q+1-N_F}{\rho}=-\sum_{i=0}^7m_i\phi_i .
 \label{eq:fermat-phase-total}
\end{equation}
There are eighteen repeated-entry contributions, three for each
of six types, and twelve residue-triple contributions, six for
each of two types. The formal proof consumes the actual
Gauss--Jacobi identity on every admitted pair and checks the
complete finite index classification. It assumes no
character-to-source identification contract.

\paragraph{A single complete cubic observation.}
Let $S(a)$ be the full source projection, retaining every
principal coordinate, and put
\[
 {\rm PC}=P_0^3+
 6P_0\bigl(2t^4X^3Y^6+t^6X^{10}Y^6\bigr),\qquad
 Q_{\rm fin}(a)=1+{\rm PC}-S(a)^3.
\]
The principal-pair correction is essential: the full cubic has
nineteen principal resonant ordered grade triples, not only $P_0^3$.
The completed cubic-spectrum theorem gives
\[
 \langle Q_{\rm fin}\rangle=1-\Phi(t,g,h,j,X,Y),\qquad
 \Phi=\sum_{i=0}^7m_i\phi_iw_i.
\]
This equality preserves the full signed spectrum and all native
weights.

For the full coherent source $S^{\rm coh}$, define
\[
 Q_{\rm coh}(a)=1+{\rm PC}-
       \bigl(\Delta S^{\rm coh}(a)\bigr)^3,\qquad
 \Delta=(1-t^2X^7)(1-t^2Y^7).
\]
When $\Delta\ne0$, compensation of both complete periods gives
$Q_{\rm coh}=Q_{\rm fin}$ pointwise. Moreover,
\[
 \langle Q_{\rm coh}\rangle-1
 =\sum_{i=0}^7m_i
       \bigl(\langle Q_i^{\rm coh}\rangle-1\bigr).
\]
This is an equality after averaging for the single full cubic
observation. It is not a product of separately averaged probes.
The proof subtracts the entire principal correction before
identifying the weighted eight-mask assembly.

\paragraph{Actual finite and joint Abel theorems.}
At any nonzero $u$, take
\[
 t=u^7,\qquad X_0=Y_0=u^{-2}.
\]
All eight native weights equal $u^{15}$. Substitution of
\eqref{eq:fermat-phase-total} proves the exact source theorem
\begin{equation}
 \langle Q_{\rm fin}\rangle
   =1+u^{15}\frac{q+1-N_F}{\rho}.
 \label{eq:fermat-finite-source}
\end{equation}
For $X=sX_0,Y=sY_0$, $0\le s<1$, both period ratios are $s^7$.
Their complete double series converge, and
$\Delta=(1-s^7)^2\ne0$. The full observation has the nontrivial
joint Abel limit
\begin{equation}
 \lim_{s\to1^-}
 \frac{\langle Q_{\rm coh}(sX_0,sY_0)\rangle-1}{u^{15}}
 =\frac{q+1-N_F}{\rho}.
 \label{eq:fermat-abel-source}
\end{equation}
To prove it, apply the eight already proved normalized Abel
limits, multiply by their exact multiplicities, and sum.
The full-source assembly holds eventually in the same radial
region, so it transfers this limit to $Q_{\rm coh}$. Lean
checks the nontrivial filter, actual convergence region and
complete source consumer.

For characteristic different from seven the Fermat septic is a
smooth genus-fifteen curve. Its Jacobi theory and Jacobian
decompositions are classical; compare
\cite[Theorems 6--7]{HidalgoRodriguez}. The new formal source
statement is the precise complete compensated observation
\eqref{eq:fermat-finite-source}--\eqref{eq:fermat-abel-source}.
Priority for that particular bridge has not been established.
Rational-prime descent, an admissible global metaplectic Fourier
realization and an independent signed analytic estimate remain
open. With $u=q^{-1/14}$ the trace extraction still costs
$q^{15/14}$; the installed QRH boundary stays $7/8$.


\section{A complete twelvefold $G_2$ coordinate calculation}
\label{sec:G2-coordinates}

This section computes the Whittaker-functional coordinates in Gao's
convention. They are not the values of a spherical Whittaker function:
that function also involves the spherical-vector expansion. A global
Fourier coefficient and a moving-twist summation formula are further
objects. The calculation below does not identify either of them.

Fix a tame unramified local field containing $\mu_{24}$, with residue
field $\mathbb F_q$, and the standard twelvefold $G_2$ quadratic datum.
Thus $q\equiv1\pmod{24}$ and the torus section sign is one. Put
$t=q^{-1/2}$ and let
\[
 G_r=q^{-1/2}\sum_{v\in\mathbb F_q^\times}\chi(v)^r\psi^{-1}(v),
 \qquad
 (g,h,j,k,e)=(G_1,G_3,G_4,G_5,G_6),
\]
where $\chi$ has exact order twelve and $\psi$ is nontrivial.
Then $|G_r|=1$ for $r\ne0\pmod{12}$,
$G_rG_{12-r}=1$, and $e^2=1$.
The coroot quadratic form and its bilinear matrix are
\[
 Q(a,b)=a^2-3ab+3b^2,\qquad
 B_Q=\begin{pmatrix}2&-3\\-3&6\end{pmatrix}.
\]
Consequently $Y_{Q,12}=12\ZZ\oplus4\ZZ$.

Normalize the free-orbit functional by $c(0,0)=1$. The source
equivariance is $c(t_0z)=c(t_0)\bar\chi(z)$, not its inverse.
The exceptional-character condition gives
\[
 c(a+12,b)=q^{-1}c(a,b),\qquad
 c(a,b+4)=q^{-1}c(a,b).
 \tag{G2.1}
\]
These signs and period directions matter. In the same convention,
Gao's root recurrence is
\[
 c(s_\alpha[y])=
 q^{\lceil\langle y,\alpha\rangle/n_\alpha\rceil-1}
 g_{\psi^{-1}}\!\left(Q(\alpha^\vee)
       (\langle y,\alpha\rangle-1)\right)^{-1}c(y).
 \tag{G2.2}
\]
Here the source $g_{\psi^{-1}}$ has modulus $q^{-1/2}$; the $G_r$
above have modulus one. See \cite[Section 3.1, Corollary 3.7,
Theorem 7.1]{Gao}. We use the free-orbit-supported functional supplied
there; the representation-theoretic identification is a cited input,
not a theorem formalized in the accompanying arithmetic modules.

\paragraph{Proposition (complete signed rectangle).}
In $0\le a<12$, $0\le b<4$, the nonzero coordinates are exactly
\[
\begin{array}{c|c|r@{\qquad}c|c|r}
(a,b)&c(a,b)&Q\bmod12&(a,b)&c(a,b)&Q\bmod12\\ \hline
(0,0)&1&0&(1,0)&tg&1\\
(0,1)&th&3&(1,2)&t^2ge&7\\
(4,1)&t^2hj&7&(6,2)&t^3gke&0\\
(4,0)&tj&4&(6,1)&t^2gke/h&9\\
(9,0)&t^2jk&9&(10,1)&t^3gkej/h&1\\
(9,2)&t^3jke&3&(10,2)&t^4gkej&4 .
\end{array}
\]
The remaining thirty-six coordinates vanish.
Writing $A=gke$, their signed numerator is
\[
\begin{split}
 N(X,Y)={}&(1+tjX^4)(1+thY)(1+t^2A h^{-1}X^6Y)\\
 &+tX(g+tjkX^8)(1+teY^2).
\end{split}
 \tag{G2.3}
\]
For every pair of nonnegative integers the period extension is unique,
and its full squared-coordinate mass is
\[
 \sum_{a,b\ge0}|c(a,b)|^2
 =\frac{1+3x+4x^2+3x^3+x^4}{(1-x^2)^2}
 =\frac{1-x^3}{(1-x)^3},\qquad x=q^{-1}.
 \tag{G2.4}
\]

\textit{Proof.}
The two shifted reflections and the pairings are
\[
 s_1[a,b]=(1-a+3b,b),\quad s_2[a,b]=(a,1+a-b),
 \qquad 2a-3b,\ -a+2b.
\]
Their effective root orders are twelve and four.
On reduction back to the rectangle, the change in the exponent of $t$ is
\[
 1-2\left\lceil\frac{\langle y,\alpha\rangle}{n_\alpha}\right\rceil
 -2\left(\left\lfloor\frac{a'}{12}\right\rfloor+
          \left\lfloor\frac{b'}4\right\rfloor\right),
\]
where $(a',b')=s_\alpha[y]$. The Gauss multiplier is
$G_{Q(\alpha^\vee)(1-\langle y,\alpha\rangle)}$.
Applying these rules from $(0,0)$ gives the displayed twelve states.
Every edge on these states obeys the same rules, including both
length-six routes to the last state; opposite Gauss factors and $e^2=1$
are retained. Explicit alternating reflection paths from $(0,0)$
reach every displayed state, proving recurrence uniqueness there.
The arithmetic checks exhaust all $48$ states and both simple roots;
this finite reduction is exact and is not a sample of primes.

Expanding (G2.3) gives precisely the twelve displayed monomials.
Every natural pair has the unique decomposition
$(a_0+12m,b_0+4n)$ in the rectangle and $m,n\ge0$.
The period factor in its squared norm is $x^{2m+2n}$.
The twelve base degrees have multiplicities $1,3,4,3,1$.
Two convergent geometric series therefore give (G2.4);
its simplification uses
$1+3x+4x^2+3x^3+x^4=(1+x)^2(1+x+x^2)$.
\hfill$\square$

\paragraph{Theorem (six coherent channels and the complete local average).}
Put
\[
 B=\frac{kej}{h},\qquad C=\frac{hj}{ge},\qquad
 D(t)=1+3t^2+4t^4+3t^6+t^8.
\]
For a twelfth root of unity $z$, rescaling the additive character
gives the complete rectangle signal
\[
 T(z)=\sum_{\substack{0\le a<12\\0\le b<4}}c(a,b)z^{Q(a,b)}.
\]
The sign choice is absorbed in $z$; as the scaling unit varies, $z$
takes every twelfth root equally often. Its exact factorization is
\[
\begin{split}
 T(z)={}&(1+t^3A)(1+tjz^4)
 +t(gz+hz^3)(1+t^2B)\\
 &+t^2(gez^7+A h^{-1}z^9)(1+C).
\end{split}
 \tag{G2.5}
\]
In particular,
\[
\begin{split}
 \frac1{12}\sum_{z^{12}=1}|T(z)|^2
 ={}&(1+t^2)|1+t^3A|^2
      +2t^2|1+t^2B|^2+2t^4|1+C|^2\\
 ={}&D(t)+2t^3(1+t^2)\Re A+4t^4(\Re B+\Re C).
\end{split}
 \tag{G2.6}
\]
These identities also hold for a formal amplitude $t\in\mathbb R$
with the unit Gauss phases fixed.

For the actual finite-field Gauss phases one has
\[
 BC=1,\qquad C^2=\chi^{-3}(3),\qquad C^4=1,
 \qquad A=\frac{J(\chi,\chi^5)}{\sqrt q}.
 \tag{G2.7}
\]
Moreover $J(\chi,\chi^5)=a_q+ib_q\in\ZZ[i]$ and
$a_q^2+b_q^2=q$. Thus for some $\epsilon_q\in\{-1,0,1\}$
the complete average at $t=q^{-1/2}$ is
\[
 1+\frac3q+\frac{4+2a_q+8\epsilon_q}{q^2}
       +\frac{3+2a_q}{q^3}+\frac1{q^4}.
 \tag{G2.8}
\]

\textit{Proof.}
The Gauss weight of each coordinate is $Q(a,b)\pmod{12}$:
the exponents of $g,h,j,k,e$ carry weights $1,3,4,5,6$.
Exactly the classes $0,1,3,4,7,9$ occur, twice each.
Grouping the twelve terms gives (G2.5). Character orthogonality
removes precisely the pairs in different classes.
The six coefficients, in the same order, are
\[
 1+t^3A,\quad tg(1+t^2B),\quad th(1+t^2B),\quad
 tj(1+t^3A),\quad t^2ge(1+C),\quad t^2A h^{-1}(1+C).
\]
Taking their squared norms proves both lines of (G2.6).
No same-class cross term has been discarded.

For completeness, the Gauss reduction in (G2.7) uses only the
duplication and triplication instances of Hasse--Davenport
\cite[(1.5)--(1.6)]{EvansGreene}. In normalized notation they give
\[
 G_1/G_5=\chi(4)^{-1}G_2e,\qquad
 G_2/G_4=\chi(4)^{-2}G_4e,\qquad
 G_1G_5/G_3=\chi^{-3}(3)G_3.
\]
Since $q\equiv1\pmod{24}$, the quadratic characters of $2$ and $3$
are one. This follows first in the prime field from the quadratic
supplementary laws, and in extensions by the norm/lift formula.
Hence $\chi(4)^3=1$, $\chi^{-3}(3)^2=1$ and $e^2=1$.
It follows that
\[
 g=kj^2,\qquad gk=\chi^{-3}(3)h^2.
\]
Substitution in the definitions of $B,C$ proves $BC=1$ and
$C^2=\chi^{-3}(3)$. Since $|C|=1$, $\Re B=\Re C$,
and $C^4=1$ gives $\Re C\in\{-1,0,1\}$.

The Gauss--Jacobi identity and $G_6^2=1$ give
$gke=J(\chi,\chi^5)/\sqrt q$.
This Jacobi sum is an algebraic integer in $\QQ(\zeta_{12})$.
The automorphism $\zeta_{12}\mapsto\zeta_{12}^5$ swaps its two
characters, so symmetry of the Jacobi sum fixes it.
Its fixed field is $\QQ(i)$, whose integers are $\ZZ[i]$.
The Gauss norm formula gives norm $q$. Expanding (G2.6) now yields
(G2.8).
\hfill$\square$

\paragraph{What the cancellation says.}
The average has no interference of order $q^{-1/2}$ or $q^{-1}$:
the first non-diagonal term is $O(q^{-3/2})$.
This is a local cancellation statement with an exact retained correction,
not a signed global prime estimate. The first correction is a Gaussian
Jacobi sum, and the next phase has only four possibilities.
Thus a twelvefold cover does not supply twelve independent arithmetic
channels in this observable.

The all-coordinate coherent sum has additional period interference.
Its generating function is
\[
 \frac{N(X,Y)}{(1-t^2X^{12})(1-t^2Y^4)}.
\]
At $X=Y=1$ its averaged squared norm is (G2.6) divided by
$(1-t^2)^4$. This differs from the coordinate mass (G2.4);
neither denominator may be used for the other quantity.

\paragraph{Formalization and priority scope.}
The maintained \texttt{G2TwelveLocal} module checks the finite support,
every signed root edge, the quadratic grading, both numerator
factorizations and the convergent mass on all natural coordinate pairs.
The accompanying \texttt{G2TwelveFourier} module treats the complete
finite Fourier average and its signed expansion. Its algebraic
Hasse--Davenport consumers retain $g=kj^2$, $gk=uh^2$, $e^2=u^2=1$
as explicit hypotheses. The finite-field Hasse--Davenport theorem,
the Jacobi fixed-field argument and the representation-theoretic
identification are cited ordinary inputs, not claimed Lean theorems.
No global Fourier coefficient, moving-twist reflection or QRH exponent
improvement follows from these arithmetic modules alone.

The explicit calculation is distinct from the prior quintic
$\GL_4$ programme. Gao supplies the source recurrence and the
distinguished-model classification, while the Gauss product identities
are classical. We have not established literature priority for the
combined coordinate and average formulas; their role here is to
replace an unexplored candidate by a completely specified local model.

\section{The quartic vector route now has a concrete comparison to inspect}

A November 3 research manuscript reflects a complete 24-cusp quartic
Gauss observation, including shared primes and dyadic masks.
It states, for $B\ge C\ge2$,
\[
 I(C,B)=K^*B^{3/4}\widehat\Phi_0(3/4)S_C+
 O_{A,\eps}\big((CB)^\eps C^{1/2}B^{3/4}(C/B)^A\big),
 \tag{10.1}
\]
where
\[
 S_C=\sum_{\substack{c\ {\rm odd\ primary,\ squarefree}\\c\ne1}}
 \mu(c)\eta(c)\prod_{\mathfrak p\mid c}(1-q_{\mathfrak p}^{-1})
 (\N c)^{-1/4}\Phi(\N c/C).
\]
This is the author's theorem statement, not a proof audited here.
It expressly gives no new zero-free region \cite{QuarticObservation}.
The importance for this map is that full vector reflection is a real
analytic object to inspect, rather than a reason to guess scalar
quartic theta coefficients.

One can calculate its proposed arithmetic consumer without assuming
that statement has been locally verified. For fixed unitary $\eta$,
the underlying Dirichlet series initially has Euler product
\[
 \sum_c\frac{\mu(c)\eta(c)\prod_{\mathfrak p\mid c}
 (1-q_{\mathfrak p}^{-1})}{(\N c)^s}
 =\frac{H_\eta(s)}{L^{(S)}(s,\eta)},
\]
\[
 H_\eta(s)=\prod_{\mathfrak p\notin S}
 \frac{1-\eta(\mathfrak p)(1-q_{\mathfrak p}^{-1})q_{\mathfrak p}^{-s}}
      {1-\eta(\mathfrak p)q_{\mathfrak p}^{-s}} .
 \tag{10.2}
\]
Omitting $c=1$ changes the series by an entire constant. For
$\Re s>0$, each denominator is nonzero, the numerator is nonzero,
and the defect is $O(q_{\mathfrak p}^{-1-\Re s})$ on compact subregions.
Thus $H_\eta$ is analytic and nonzero there. A zero of the denominator
$L$ is not erased by this correction.

At balanced scales $B=C=X$, a zero at real part $\beta$ gives the
Mellin pole exponent
\[
 X^{3/4}X^{\beta-1/4}=X^{\beta+1/2}.
\]
This is a pole/oscillation scale, not a lower bound at every cutoff;
the smoothing must not vanish at the relevant pole.
The claimed comparison error is $X^{5/4+\eps}$.
Therefore an \emph{independently proved} bound
\[
 |I(X,X)|\ll_{\eta,\eps}X^{5/4+\eps}
 \tag{10.3}
\]
for this actual signed observation, with the needed compatible analytic
hypotheses, would have the correct exponents to exclude $\beta>3/4$.

The absolute envelope for $S_X$ is only $X^{3/4+\eps}$, giving a
main envelope $X^{3/2+\eps}$ for $I$.
Equation (10.1) alone does not improve it. Its error can be much smaller
at $B>C$ without providing any bound on $S_C$.
An estimate for every arbitrary bounded first coefficient cannot give
the desired saving: replacing the signed coefficient by a positive one
can saturate this main envelope. The required arithmetic restriction
is the actual Möbius--Hecke coefficient, not arbitrary bounded data.

Thus the quartic comparison has useful precision if its proof checks;
it is not already an independent cancellation theorem. This route must
be compared with the project's older Gaussian packets at the level
of their exact sums before any old bound is transferred.

\section{How to explore genuinely different mechanisms}

Three distinctions prevent the investigation collapsing into repeated
versions of the same failed estimate.

\textbf{Retain a vector when the scalar coefficient is not closed.}
For a finite cusp vector $\mathbf T$ with reflection matrix $M$,
start with $\mathbf T(s)=M(s)\mathbf T(1-s)$.
A scalar row must be followed through every component it generates.
Discarding an unknown component is not a projection theorem.
If a proposed linear compensation annihilates a residue subspace $E$,
it preserves the target signal only if that signal has a nonzero
image in the corresponding quotient. If the signal lies in $E$,
that compensation necessarily kills it.
Fixed-level finite dimension also gives no uniform bound when the
moving conductor enlarges the local data.

\textbf{Use multiple root equations as multiple equations.}
Weyl group multiple Dirichlet series provide meromorphic continuation
and functional equations in several variables \cite{CG}.
Metaplectic Eisenstein Whittaker coefficients admit uniform-degree
crystal or Gelfand--Tsetlin descriptions in type $A$ \cite{BBF}.
The opportunity is to retain neighboring-coordinate coupling when
transforming and averaging. Evaluating the extra variables or bounding
their sums before using the other functional equations can lose this
information. Neither theorem supplies a moment saving for free.

\textbf{Compare the effective order and conductor, not the cover's name.}
A transform of an order-eight family can descend through quartic and
quadratic characters. An order-ten family cycles through fifth-order
information. The root-number coupling and conductor of a powered index
may differ from its norm. An actual route must write these exact
couplings, their zero masks and the global profile costs before applying
the appropriate character estimate.

A complete candidate receipt should answer:
\begin{enumerate}
\item Which representation, cover, field and global Fourier coefficient?
What is the nonzero normalization?
\item Which exact local generating series and coefficient index?
Does its global product really yield a reciprocal target $L$-factor?
\item What is the full moving-twist reflection, including auxiliary
coordinates, small primes, units and archimedean profiles?
\item What are the actual row and column lengths, conductor, height and
exceptional families after every transform?
\item Which \emph{independent} inequality bounds the complete signed
observable, and what boundary follows after all errors are included?
\end{enumerate}
This is a screening framework, not a promise that five routine steps
finish a new zero-free theorem.

\section{A sharp inequality for the Paley phase sector}

The full Fermat observation supplies a concrete stability question that
is different from recovering a phase by a classical bispectrum.
Write \(\Phi(g,h,j)=\texttt{phaseTotal}(g,h,j)\), with all eight
multiplicities from the preceding construction. In this section we
restrict explicitly to
\[
 |g|=1,\qquad h=g,\qquad j=g^{-1}.
\]
This is an additional sector hypothesis. The previously proved diagonal
Jacobi Frobenius theorem alone does not discharge it for arbitrary Gauss
inputs. We do not claim that it covers every arithmetic source.

Put \(c=\Re g\), \(s=|\Im g|\), \(t=1/\sqrt2\), \(v=\sqrt7\), and define
the explicit sector defect
\[
 d_P(g)=\min\{6-6c,\;6+c-vs\}.
\]
The statement below uses this formula as its definition. An identification
with a Euclidean distance to all seven coherent phase triples is not
among the formal claims of this section.

\paragraph{Theorem (sharp full-source inequality on the Paley sector).}
For every unit complex number \(g\),
\[
 30-\Re\Phi(g,g,g^{-1})
 \ \ge\left(6+\frac3{\sqrt2}\right)d_P(g).
\]
The constant is optimal uniformly throughout this sector. It is attained
at the noncoherent point
\[
 g=-\frac1{\sqrt8}+i\sqrt{\frac78}.
\]
\paragraph{Proof.}
Expanding the actual weighted eight-phase table, rather than discarding
either cubic sector, gives
\[
 \Phi(g,g,g^{-1})
 =6(g+g^{-1})^3-9(g+g^{-1})=48c^3-18c.
\]
Thus it suffices to prove, for every \(c^2+s^2=1\), \(s\ge0\),
\[
 (6+3t)\min\{6-6c,6+c-vs\}\le30+18c-48c^3.                 \tag{P1}
\]
Here \(t>0\), \(t^2=1/2\), \(v\ge0\), and \(v^2=7\).
If \(c\ge t/2\), subtraction of \((6+3t)(6-6c)\) factors as
\[
 6(1-c)(8c^2+8c-1-3t)\ge0.
\]
Indeed \(c\le1\), \(8c^2\ge1\), and \(8c\ge4t\).

Suppose \(c\le t/2\), and retain the signed expression
\[
 A(c)=6+18t+(-12+3t)c+48c^3.
\]
We need \(A(c)\le(6+3t)vs\). If \(c\le-7/8\), let
\(w=-7/8-c\ge0\). The exact identity
\[
 A(c)=-\frac{501}{32}+\frac{123}8t
 -\left[\left(\frac{393}4+3t\right)w+126w^2+48w^3\right]
\]
proves \(A(c)\le0\); for example \(7/10\le t\le71/100\)
follows directly from \(t^2=1/2\), \(t>0\).
This tail is treated before squaring.

For \(-7/8\le c\le t/2\), the exact residual is
\[
 7(6+3t)^2(1-c^2)-A(c)^2
   =144(c+t/2)^2Q_t(c),                                  \tag{P2}
\]
where
\[
 Q_t(c)=-16c^4+16tc^3+(2-2t)c^2-(3+16t)c+\frac{19}4+2t.
\]
For \(c=-7/8+(7/8+t/2)z\), \(0\le z\le1\), the complete
positive certificate is
\[
 Q_t(c)=\sum_{r=0}^4 \binom{4}{r}b_rz^r(1-z)^{4-r},
\]
with
\[
\begin{aligned}
 b_0&=-121/256+15t/4,& b_1&=9+1741t/128,\\
 b_2&=893/192+1247t/192,& b_3&=23/8+27t/8,\\
 b_4&=5/4+t/4.
\end{aligned}
\]
Every \(b_r\) is positive. Hence (P2) is nonnegative.
Since \(s^2=1-c^2\) and \((6+3t)vs\ge0\), it follows that
\(A(c)\le(6+3t)vs\), proving (P1).

At \(c=-t/2\), \(s=vt/2\), the circle condition holds and
\[
 d_P=6-4t>0,\qquad
 30+18c-48c^3=30-6t=(6+3t)(6-4t).
\]
This proves equality and rules out any larger uniform constant.
\hfill$\square$

All the polynomial identities, the signed tail, the uniform circle
inequality, its symbolic attainer, optimality, and the full-source
consumer are kernel checked in
\texttt{G2SevenPaleyStability.lean}. In particular the statement is
uniform, not an inference from a numerical scan.

The unrestricted three-phase extension at the sharp sector constant
remains a conjecture. A finite
diagnostic on 132768 grid and random triples, followed by numerical
optimization, found the same candidate constant; this is neither a proof
nor a global QRH estimate. Two low-degree polynomial certificate attempts
did not close and are recorded as unsuccessful. Plain cyclic bispectrum
reconstruction is classical, and priority for this particular sharp
source inequality has not been established. The installed zero-free
boundary remains \(7/8\).

\subsection{Arithmetic discharge of the symmetry sector}

The preceding sharp inequality has an actual finite-field consumer.
Let \(F\) be any finite field of prime characteristic \(p\), let
\(\chi:F\to\mathbb C\) have exact order seven, and let \(\psi\) be
primitive. In this subsection the additive character satisfies the
explicit condition
\[
 \psi(x^p)=\psi(x)\qquad(x\in F).                         \tag{P3}
\]
Choose a positive real \(\rho\) with \(\rho^2=q=\#F\), and write
\(G_n=G(\chi^n,\psi)/\rho\).
The condition concerns Frobenius on the source field, not a
characteristic-\(p\) operation on the complex target.

\paragraph{Theorem (sharp signed Fermat bound in the order-three classes).}
If \(p\equiv2\) or \(4\pmod7\), then \(G_2=G_1\),
\(G_3=G_1^{-1}\), \(|G_1|=1\), and
\[
 \Re\frac{q+1-N_{\mathrm{chart}}}{\rho}
 \ \ge -30+\left(6+\frac3{\sqrt2}\right)d_P(G_1).          \tag{P4}
\]
The count includes the complete infinity chart and all thirty
nondegenerate Jacobi terms.

\paragraph{Proof.}
Reindexing the entire Gauss sum by the finite-field automorphism
\(x\mapsto x^p\), using (P3), gives
\[
 G(\chi^p,\psi)=G(\chi,\psi),\qquad G_{np}=G_n.
\]
Together with \(G_{n+7}=G_n\), the classes \(p\equiv2,4\)
give \(G_2=G_4=G_1\). The previously proved complement identity
\(G_4=G_3^{-1}\) gives \(G_3=G_1^{-1}\).
The exact conjugate product is
\[
 G(\chi,\psi)\,\overline{G(\chi,\psi)}=q.
\]
Since \(\rho\) is real and positive, division by \(\rho^2=q\)
proves \((\Re G_1)^2+(\Im G_1)^2=1\).
Finally, the already established full Fermat readout has the sign
\[
 T:=\frac{q+1-N_{\mathrm{chart}}}{\rho}
   =-\Phi(G_1,G_2,G_3).
\]
Substitute the proved sector identities into the sharp source inequality:
\(30+T\ge(6+3/\sqrt2)d_P(G_1)\), which is (P4).
\hfill$\square$

These steps, including the unit normalization and the final signed
point-count consumer, are kernel checked in
\texttt{G2SevenFrobeniusStability.lean}. The additive invariance (P3)
is retained in the theorem signature; it is not replaced by an assumption
that all primitive characters are invariant. Characteristic two is
included: the divisions by two occur in the real/complex comparison,
not in \(F\).

This arithmetic transfer resolves a symmetry-sector applicability question.
It does not control the split class \(p\equiv1\pmod7\), discharge
rational-prime descent, or estimate a global signed moment. In particular
(P4) is a source-coherence refinement of a point-count bound; it supplies
no new global zero-free exponent by itself.

\subsection{The scalar region and its exact limitation}

A supporting calculation for the unrestricted phase problem can be
settled completely. Let \(F_0,\ldots,F_6\) be arbitrary real numbers with
\[
 \sum_iF_i=7,\qquad \sum_iF_i^2=49,\qquad F_i\le M,\qquad M>1.
\]
The exact weighted-square identity is
\[
\begin{split}
 &\sum_i(M-F_i)\big((M-1)F_i-(M-7)\big)^2\\
 &\quad=(M-1)\left[(M-1)\left(343-\sum_iF_i^3\right)
                     -42M(7-M)\right].
\end{split}
\]
Consequently
\[
 (M-1)\left(343-\sum_iF_i^3\right)\ge42M(7-M).             \tag{P5}
\]
Put \(t=1/\sqrt2\). Throughout the complete range
\(1<M\le1+4t=1+2\sqrt2\), the additional signed slack is
\[
 42M(7-M)-7(6+3t)(M-1)(7-M)
   =21t(7-M)(1+4t-M)\ge0.
\]
Together with (P5), this proves
\[
 \frac{343-\sum_iF_i^3}{7}\ge(6+3t)(7-M).                 \tag{P6}
\]
These are kernel-checked statements about every real seven-vector with
the stated two moments. They are supporting scalar input; a formal
source-to-real-vector or nearest-phase-distance consumer is not asserted
by this module.

The limitation of the two-moment relaxation is also exact. Take
\(v=\sqrt{15}\), two entries \(1+v\), and five entries \(1-2v/5\).
This vector satisfies both moments and the upper bound \(M=1+v\), but
\[
 \frac{343-\sum_iF_i^3}{7}=30-\frac{18}5v
 <(6+3t)(6-v).
\]
The positive difference is
\((6-v)(3t-2v/5)\). With the two high entries adjacent, its first cyclic
correlation is
\[
 \sum_{i\in\mathbb Z/7\mathbb Z}F_iF_{i+1}=\frac{98}5,
\]
rather than zero. Thus this counterexample excludes an extension using
only those moments; it does not refute the full cyclic phase problem.
All these scalar statements, including the explicit counterexample and
its failed cyclic correlation, compile in
\texttt{G2SevenScalarStability.lean}. The weighted-square envelope is a
supporting classical calculation, not an originality or QRH claim.

At the proposed sharp constant, the remaining phase-geometry range is
\(1+2\sqrt2<M<7\), after the
Fourier moment normalization is supplied. The attempted cubic and
quartic positive-polynomial certificates have not settled that range.
Their numerical coefficient fits and projected positive matrices are
diagnostics, not formal theorems or infeasibility proofs.

\section{A uniform cyclic stability theorem}

The earlier sharp sector calculation has a weaker uniform extension
on the complete cyclic correlation domain. This extension uses all
correlations, not merely the first two moments.

\paragraph{Statement.}
Let \(y_0,\ldots,y_5\) be arbitrary real numbers, put
\[
 m=1-\sum_{r=0}^5y_r,\qquad
 x=(m,y_0,\ldots,y_5),\qquad
 C_h(x)=\sum_{i\in\ZZ/7\ZZ}x_ix_{i+h},
\]
and assume \(y_r\le m\) for every \(r\) and
\(C_1(x)=C_2(x)=C_3(x)=0\). Then
\begin{equation}\label{eq:uniform-eight}
 7\left(1-\sum_{i=0}^6x_i^3\right)\ge8(1-m).
\end{equation}
The coordinates need not be nonnegative. There is no sampling,
Frobenius-sector or asymptotic hypothesis. The compiled declaration is
\texttt{Metaplectic.G2Seven.Uniform.}\allowbreak
\texttt{normalized\_stability}.

\paragraph{Proof by an exact signed certificate.}
First, the correlation equations, unit mass and peak condition imply
\[
 \sum_i x_i^2=1,\qquad 2\le7m,\qquad m\le1.
\]
The six multiplicative automorphisms of \(\ZZ/7\ZZ\) fix coordinate zero,
preserve its dominance, and carry \(C_1\) into the three assumed
correlations or their reverses. Thus every transported first
correlation is exactly zero. These facts are proved in
\texttt{G2SevenCyclicGeometry.lean}.

Set
\[
 T(y)=7\left(1-\sum_i x_i^3\right)-8(1-m).
\]
The complete certificate is specified by the integer data in
\texttt{G2SevenUniformCertificateData.lean}. It consists of
\(79\) source-defined polynomial columns \(P_j\), integer weights
\(w_j\), an explicit positive integer \(D\), and the exact identity
\begin{equation}\label{eq:uniform-eight-certificate}
 D\cdot704643072\,T(y)=\sum_{j=0}^{78}w_jP_j(y).
\end{equation}
The data file gives every coefficient and weight; they are not
floating-point approximations. Its first \(39\) weights are
nonnegative. The remaining \(40\) weights retain their signs.

Here is the structure that makes the identity a proof.
The first \(38\) columns are positive constant multiples of sums of
squares over all six multiplicative automorphisms. Their extra factors
are respectively \(m-y_r\), \(1-m\), or \(7m-2\), which are nonnegative.
Column \(38\) is a positive constant multiple of
\[
 (1-m)\bigl[
 (m-y_0)(m-y_1)(m-y_3)
 +(m-y_2)(m-y_4)(m-y_5)\bigr],
\]
and is nonnegative too. Each of the final \(40\) columns is a sum
of a transported first correlation times its prescribed monomial.
Every such column vanishes; no estimate by its absolute value is used.

For completeness, identity~\eqref{eq:uniform-eight-certificate}
is checked coefficient by coefficient in the \(80\) cyclic monomial
orbits of the degree-at-most-five polynomial space used by the
certificate. The \(79\) individual polynomial expansions are proved
by exact ring identities in the sixteen bounded modules
\texttt{G2SevenUniformCertificateBlock00.lean} through
\texttt{Block15.lean}. The columns module assembles these identities.
The integer coefficient identities and the nonnegative-weight check
then give~\eqref{eq:uniform-eight-certificate} over the reals.
Consequently its right side is nonnegative. Division by the positive
factor \(704643072D\) proves \(T(y)\ge0\), which is
\eqref{eq:uniform-eight}. The complete signed assembly and final
inequality compile in \texttt{G2SevenUniformStability.lean}.
\(\square\)

\paragraph{One complete column record.}
The orbit basis is unambiguous: take distinct monomial orbits under
multiplication on the nonzero residues modulo \(7\), sum each orbit
without repetitions, and order representatives first by increasing
degree and then by decreasing lexicographic exponent tuple
\((a_0,\ldots,a_5)\). There are \(1,1,4,10,22,42\) such orbits in
degrees \(0,1,2,3,4,5\), respectively.

For column \(j=38\), the normalization is \(s=704643072\), and
\(P_{38}/s\) is exactly the bracketed Fano polynomial above, including
its factor \(1-m\). Its coefficient vector divided by \(s\), in this
ordered basis, starts with
\[
 \begin{gathered}
0,2,-7,-14,-14,-14,8,25,24,25\\
24,24,49,51,49,49,-3,-13,-12,-13\\
-12,-12,-20,-38,-42,-38,-38,-18,-38,-37\\
-37,-38,-38,-37,-18,-79,-76,-76,0,0
 \end{gathered}
\]
at indices \(0,\ldots,37\), and every coefficient at indices
\(38,\ldots,79\) is zero. Its integer weight \(w_{38}\) is obtained
by concatenating the following decimal lines:
\begin{quote}\small\ttfamily
286105408823985556996573140820681364628834890004\\
906256298699824475477629399121996992813147441358\\
375819839913511439612448084272302987244478252721\\
126090830377190399655865404928538494634294919766\\
129056694624404262727450490774991051960175599889\\
300983338219247017915030777948228508416
\end{quote}
Thus the record specifies its index, normalization, full polynomial,
all \(80\) coefficients and exact weight. Nonnegativity follows from
the four displayed nonnegative factors and \(w_{38}\ge0\); kernel
ring identities verify the expansion. The data module contains all
\(79\) records and their common positive normalizer, as part of the
same maintained source tree. Public release links for this new
certificate remain pending publication.

This theorem covers the entire feasible peak range at constant \(8\).
It does not establish the sharper candidate \(6+3/\sqrt2\) on that
range. A separately compiled Fourier or matrix-distance transport,
arithmetic membership statement and global automorphic estimate are
not part of this declaration. In particular the QRH boundary remains
\(7/8\).

\subsection{A concrete signed arithmetic edge candidate}

There is an arithmetic question beyond the uniform phase geometry.
For a rational prime \(p\equiv1\pmod7\), let \(a,b,c,d\) denote the
unnormalized Jacobi sums \(J_{1,2},J_{1,1},J_{1,3},J_{2,2}\), and
write \(A=a/\sqrt p\), \(B=b/\sqrt p\), \(C=c/\sqrt p\), \(D_0=d/\sqrt p\).
The already proved full Fermat phase formula retains
\[
 \Phi=6\Re(B+C+D_0)+12\Re(A),\qquad A=BD_0/C.
\]
The signed algebra behind a possible stronger estimate is
\[
 (b-c)(d-\overline b)(\overline c-\overline d)
   =p\bigl[(a+c-b-d)-\overline{(a+c-b-d)}\bigr].
\]
It follows by an eight-term expansion from
\(b\overline b=c\overline c=d\overline d=p\),
\(ac=bd\), and its conjugate. The generic commutative-ring identity
and the supporting inequality
\(27uvw\le(u+v+w)^3\) for \(u,v,w\ge0\) are maintained in
\texttt{G2SevenJacobiEdgeAlgebra.lean}; neither is claimed as a new
classical inequality.

\paragraph{Theorem (complete unit-phase edge certificate).}
Let \(B,C,D_0\) be arbitrary complex numbers of modulus one and put
\(A=BD_0\overline C\). Define
\[
 \Delta=30-6\Re(B+C+D_0)-12\Re A,\quad
 (e_1,e_2,e_3)=(|1-BD_0|^2,|B-C|^2,|C-D_0|^2).
\]
Then
\[
 \Delta\ge2(e_1+e_2+e_3),\qquad
 \Delta^3\ge216e_1e_2e_3.
\]
Consequently, for any real \(p\ge0\), the explicit arithmetic premise
\(pe_1e_2e_3\ge7\) implies
\[
 p\bigl(30-|6\Re(B+C+D_0)+12\Re A|\bigr)^3\ge1512.
\]
No angular-sector hypothesis is required.
\emph{Proof.}
Set
\[
 U_1=\overline{BD_0},\qquad U_2=D_0\overline C,\qquad
 U_3=C\overline B,
\]
and retain the complete six-square packet
\[
\begin{aligned}
 L_1&=2(1+U_1)-U_1(A+B+C+D_0),\\
 L_2&=2(1+U_2)-(A+D_0+\overline B+\overline C),\\
 L_3&=2(1+U_3)-(\overline A+\overline B+C+D_0),\\
 R_1&=(1+U_1)(1-U_2),\\
 R_2&=(1+U_2)(1-U_3),\\
 R_3&=(1+U_3)(1-U_1).
\end{aligned}
\]
The exact signed identity is
\[
 4\bigl[\Delta-2(e_1+e_2+e_3)\bigr]
 =\sum_{j=1}^3|L_j|^2+\sum_{j=1}^3|R_j|^2
   +12|1-A|^2.
\]
Its proof is a polynomial identity modulo
\(B\overline B=C\overline C=D_0\overline{D_0}=1\);
the maintained commutative-ring proof has an exact signed ideal
certificate, and its complex specialization checks every conjugation.
Nonnegativity gives the first inequality. The three-variable product
inequality then gives
\(\Delta^3\ge8(e_1+e_2+e_3)^3\ge216e_1e_2e_3\).
Negating all three phases also negates \(A\) and the complete trace,
while preserving every \(e_j\). Applying the same argument at both
edges and multiplying by \(p\ge0\) proves the last assertion.
All five declarations, including this absolute-value consumer, are
maintained in \texttt{G2SevenJacobiUnitCircle.lean}.\hfill\(\square\)

\paragraph{A common counting profile for every positive power (proved).}
A mathematical specialization of the maintained general semiring theorem
\texttt{FiniteRootPower.exists\_counting\_profile} is the following.
Let \(n\ge0\), let \(S\) be a finite subset of an arbitrary index set,
and let a complex-valued function \(f\) take values in
\(\{0,z^0,\ldots,z^{n-1}\}\). There are natural numbers \(v_j\)
such that the same vector satisfies
\[
 \sum_{j=0}^{n-1}v_j=\#\{x\in S:f(x)\ne0\},\qquad
 \sum_{x\in S}f(x)^r=\sum_{j=0}^{n-1}v_jz^{jr}
 \quad\text{for every integer }r\ge1.
\]
The case \(n=0\) is included, and no primitive-root or distinctness
hypothesis is needed. The maintained declaration works over every
commutative semiring with decidable equality, rather than only
over the complex numbers.

For a proof, add the elements of \(S\) one at a time. A zero value
changes neither the nonzero count nor any positive-power sum.
For a nonzero value choose one of its prescribed power labels and
increment that coefficient. The same label works for all \(r\ge1\).
This is a standard finite counting argument. Its role here is to
preserve one common integer coefficient vector across all character
powers; it supplies no analytic character-sum estimate.

\paragraph{Actual simultaneous Jacobi twists, including the principal row (proved).}
Let \(F\) be any finite field of cardinality \(q\), let \(n>0\),
let \(z\) be a primitive \(n\)-th root, and let \(\chi,\psi\)
be multiplicative characters with \(\chi^n=\psi^n=1\).
The maintained theorem
\texttt{Metaplectic.JacobiPower.actual\_counting\_profile}
provides one vector \(v\in\mathbb N^n\) satisfying
\[
 \sum_jv_j=q-2,\qquad
 J(\chi^r,\psi^r)=\sum_{j=0}^{n-1}v_jz^{jr}
 \quad\text{for every integer }r\ge0.
\]
For \(r>0\), apply the counting theorem to
\(f(x)=\chi(x)\psi(1-x)\). Its nonzero support is exactly
\(F\setminus\{0,1\}\), and \(f(x)^n=1\) there.
For \(r=0\), use the actual principal value
\(J(1,1)=q-2\); this avoids incorrectly raising a zero
character value to the zeroth power. The case \(q=2\) has zero
mass and is included. This applies to both diagonal and mixed
Jacobi channels, without assuming that either character has exact
order \(n\). It discharges their common counting-profile bridge;
it is not an analytic bound or an originality claim.

\paragraph{Cyclotomic nonvanishing (proved).}
For arbitrary integers \(v_0,\ldots,v_6\), a primitive seventh root
\(z\), and a rational prime \(p\), write
\(T_r=\sum_{j=0}^6v_jz^{rj}\).
If \(|T_1|^2=p\), then \(T_1\ne T_3\). If also
\(|T_2|^2=|T_4|^2=p\), then the full signed product satisfies
\[
 (T_1-T_3)(T_2-T_6)(T_4-T_5)\ne0.
\]
Both assertions are maintained and kernel checked in
\texttt{G2SevenCyclotomicNonvanishing.lean}.

Indeed, the existing primitive-root coefficient theorem says that an
integer linear combination of all seven powers vanishes precisely
when all seven coefficients agree. Apply it to \(T_1-T_3\).
Its zeroth coefficient is zero, so every coefficient difference is
zero. Multiplication by three permutes the six nonzero residues
in one cycle. Thus \(v_1=\cdots=v_6\), whence
\(T_1=v_0-v_1\) is an integer, contradicting \(|T_1|^2=p\).
Apply the same argument at the primitive roots \(z^2,z^4\);
their third powers are \(z^6,z^5\), respectively. No individual
beta factor is discarded or replaced by an unsigned average.

\paragraph{The integral quadratic anti-invariant (proved).}
Let \(z\) be primitive of order seven, and let \(v\in\mathbb Z^7\)
satisfy
\[
 v_0=0,\quad v_6=-v_1,\quad v_5=-v_2,\quad v_4=-v_3,
 \qquad T_v(z^2)=T_v(z).
\]
The maintained theorem \texttt{septic\_integral\_anti\_invariant}
proves
\[
 T_v(z)=v_1\tau,\qquad
 \tau=z+z^2+z^4-z^3-z^5-z^6,\qquad \tau^2=-7.
\]
The square identity is separately maintained as
\texttt{septic\_tau\_square}. The complete cyclotomic coefficient
relation applied to \(T_v(z)-T_v(z^2)\) forces equality of the
three quadratic-residue coefficients; the opposite-coefficient
conditions determine the remaining three. Reducing the square of
\(\tau\) by \(1+z+\cdots+z^6=0\) gives the displayed identity.
Thus \(|\tau|^2=7\), and every nonzero integral multiple obeys
\(|m\tau|^2=7m^2\ge7\), as proved by
\texttt{septic\_anti\_invariant\_norm\_gap}.
These three declarations are kernel checked in
\texttt{G2SevenJacobiAntiInvariant.lean}. Their actual Jacobi
instantiation and signed quantitative consumer are proved below.

\paragraph{Actual inverse-coordinate Jacobi symmetry (proved).}
For arbitrary complex multiplicative characters on any finite field,
including principal characters, the maintained identity is
\[
 J(\chi,\psi)=\psi(-1)J((\chi\psi)^{-1},\psi).
\]
It follows by the bijection \(x\mapsto x^{-1}\), with zero fixed,
and \(1-x^{-1}=-(1-x)x^{-1}\) for \(x\ne0\).
The zero term vanishes on both sides. The declaration
\texttt{Metaplectic.JacobiPower.inverse\_symmetry} retains
the \(\psi(-1)\) factor and is maintained in
\texttt{JacobiPowerProfile.lean}.

\paragraph{The complete coupled integral packet (proved).}
For \(u,v\in\mathbb Z^7\), put
\(U_r=\sum_{j=0}^6u_jz^{rj}\) and
\(V_r=\sum_{j=0}^6v_jz^{rj}\).
If \(U_2=U_1\) and \(U_5=U_6\), the maintained declaration
\texttt{septic\_integral\_packet} proves
\[
 (U_1+V_3-V_1-V_2)-(U_6+V_4-V_6-V_5)=M\tau,
 \qquad
 M=u_1-u_6+v_3+v_5+v_6-v_1-v_2-v_4\in\mathbb Z.
\]
Here every sum has the same support \(\{0,\ldots,6\}\).
For a direct proof, put \(S=v_3+v_5+v_6-v_1-v_2-v_4\) and take
\[
 w=(0,M,u_2-u_5+S,u_3-u_4-S,
       -(u_3-u_4-S),-(u_2-u_5+S),-M).
\]
Reduction by \(z^7=1\) identifies \(T_w(z)\) with the complete
eight-term expression above, while
\[
 T_w(z^2)=(U_2+V_6-V_2-V_4)-(U_5+V_1-V_5-V_3).
\]
The two hypotheses make these evaluations equal. The preceding
integral anti-invariant theorem gives \(T_w(z)=w_1\tau=M\tau\).
The full proof, including both finite polynomial identities, is
kernel checked in \texttt{G2SevenJacobiAntiInvariant.lean}.

\paragraph{Actual character norms, symmetries and beta nonvanishing (proved).}
Let \(F\) be a finite field of prime characteristic \(p\), let
\(q=\#F\), and let \(\chi\) have exact order seven. Write
\(B_r=J(\chi^r,\chi^r)\), \(b=B_1\), \(c=B_3\), \(d=B_2\), and
\(a=J(\chi,\chi^2)\). The maintained arithmetic packet proves
\[
 |B_1|^2=|B_2|^2=|B_3|^2=|B_4|^2=q,\qquad
 J(\chi^2,\chi^4)=a,\quad J(\chi,\chi^3)=c,\quad \chi(-1)=1.
\]
The norm proof consumes the existing Jacobi inverse-product identity
and character conjugation. The two symmetries consume the actual
inverse-coordinate identity, with their odd-order sign factors
discharged by the existing character theorem.

If additionally \(q=p\), the complete actual product satisfies
\[
 (B_1-B_3)(B_2-B_6)(B_4-B_5)\ne0.
\]
To prove this, apply the simultaneous counting profile to \((\chi,\chi)\)
and cast that one natural vector to integers. Its evaluations at
\(z^r\) are exactly the \(B_r\). The three required norm hypotheses
at \(r=1,2,4\) are now the actual norms proved above, so the
maintained cyclotomic nonvanishing theorem applies directly.
This application needs no additive character.

For a primitive additive character \(\psi\), put \(g_r=g(\chi^r,\psi)\).
The four existing Gauss--Jacobi identities give
\[
 bg_2=g_1^2,\qquad dg_4=g_2^2,\qquad
 ag_3=g_1g_2,\qquad cg_4=g_1g_3.
\]
Their exact coupled subtraction yields
\((ac-bd)g_3g_4=0\). Primitivity and the nontriviality of
\(\chi^3,\chi^4\) give \(g_3g_4\ne0\), hence \(ac=bd\).
These three actual-family declarations are kernel checked in
\texttt{G2SevenJacobiArithmetic.lean}. They are supporting arithmetic
applications; no new analytic bound or priority claim is made here.

\paragraph{Complete actual arithmetic phase edge (proved).}
Let \(F\) have rational-prime cardinality \(p\), let \(\chi\) have
exact order seven, and choose a primitive additive character.
Put
\[
 a=J(\chi,\chi^2),\quad b=J(\chi,\chi),\quad
 c=J(\chi^3,\chi^3),\quad d=J(\chi^2,\chi^2).
\]
For any \(\rho>0\) with \(\rho^2=p\), set
\(B=b/\rho,\ C=c/\rho,\ D_0=d/\rho,\ A=a/\rho\). Then
\[
 \Phi=6\Re(B+C+D_0)+12\Re A,\qquad
 \boxed{\ p(30-|\Phi|)^3\ge1512\ }.
\]
All arithmetic premises of this phase bound are discharged.

\emph{Proof.}
The actual counting profiles \(u,v:\{0,\ldots,6\}\to\mathbb Z\)
represent respectively \(J(\chi^r,\chi^{2r})\) and
\(J(\chi^r,\chi^r)\) at every natural twist \(r\).
Mixed-channel symmetry gives \(U_2=U_1\); conjugation gives
\(U_5=U_6\). The complete integral packet therefore yields
\[
 H=(a+c-b-d)-\overline{a+c-b-d}=m\tau,\qquad
 m=u_1-u_6+v_3+v_5+v_6-v_1-v_2-v_4,
\]
where \(\tau^2=-7\). No term of either profile is discarded.
The full signed algebraic identity and \(ac=bd\) give
\[
 \beta=(b-c)(d-\bar b)(\bar c-\bar d)=pH.
\]
Actual same-orbit prime-field nonvanishing gives \(\beta\ne0\),
hence \(m\ne0\) and \(|H|^2=7m^2\ge7\).
The norm identities imply \(|B|=|C|=|D_0|=1\), and
\(ac=bd,\ c\bar c=p\) imply \(A=BD_0\bar C\).
Exact positive normalization now gives
\[
 p\,|1-BD_0|^2|B-C|^2|C-D_0|^2=|H|^2\ge7.
\]
The already proved two-sided unit-phase inequality applies, giving
the displayed constant \(1512\).
The four maintained declarations are
\texttt{septic\_actual\_integral\_packet},
\texttt{septic\_actual\_beta\_identity},
\texttt{septic\_actual\_packet\_norm\_gap} and
\texttt{septic\_actual\_unit\_edge} in
\texttt{G2SevenJacobiPrimeEdge.lean}.

For the previously defined Fermat count \(N_p\), the remaining
formal scalar specialization of the full readout is
\[
 (N_p-p-1)/\rho=6\Re(B+C+D_0)+12\Re(a/\rho).
\]
With this specialization, the phase theorem gives the candidate
point-count consequence
\[
 |N_p-p-1|\le30\sqrt p-6\sqrt[3]7\,p^{1/6}.
\]
The phase theorem is complete; this direct scalar point-count
consumer has not yet been checked as a named Lean declaration.
An exact diagnostic checked the Jacobi norm identities and proposed
bound at all \(377\) primes \(p\le20000\) with \(p\equiv1\pmod7\).
The earlier Laurent residual check has been replaced by the maintained
universal ring proof and complete complex consumer above. The finite
prime computations remain falsification evidence only.
Square extension fields can reach a Hasse edge, so a uniform claim
over all finite fields is excluded. This does not show that prime
fields are the maximal valid domain; a nonsquare-cardinality extension
remains to be proved. Priority and the relation to existing
CM/Fermat point-count bounds have not been established. Even a
proof of this candidate would require a further global signed
consumer before it could give a QRH exponent improvement.

\section{Expanding the target $L$-functions is a separate problem}

There are two goals that should not be conflated: lowering the boundary
for the existing Hecke family, and reaching new families of $L$-functions.
The arbitrary-number-field extension already reaches finite-order Hecke
characters, Dedekind zeta functions and, by induction, monomial Artin
$L$-functions in its stated scope \cite[Corollary 19.2]{LiuField}.
An exceptional group in the construction does not by itself turn the
detected Hecke $L$-function into that group's standard $L$-function.

Here is a concrete coefficient test for genuinely higher-degree targets.
If an unramified target has local factor
\[
 L_{\mathfrak p}(s,\pi)=
 \prod_{j=1}^d(1-\alpha_{\mathfrak p,j}q_{\mathfrak p}^{-s})^{-1},
\]
then its inverse requires the complete polynomial
\[
 L_{\mathfrak p}(s,\pi)^{-1}
 =\sum_{r=0}^d(-1)^r
 e_r(\alpha_{\mathfrak p,1},\ldots,\alpha_{\mathfrak p,d})
 q_{\mathfrak p}^{-rs}.
\]
For $d=2$ this is
\[
 1-\lambda_\pi(\mathfrak p)q_{\mathfrak p}^{-s}
   +\omega_\pi(\mathfrak p)q_{\mathfrak p}^{-2s}.
\]
The quadratic term is essential. A construction that only produces a
squarefree Mobius coefficient times a scalar character has not produced
this inverse, even if it uses a higher-degree metaplectic cover.

Higher-rank or vector coefficients could supply the additional local
sectors. The first test is to recover every elementary-symmetric term
with its correct sign and index, with one admissible global functional.
The global transformation and the analytic estimate would then have to
include the target representation's conductor and height as well.
This algebraic test identifies what an extension to degree-two or
higher automorphic targets must add; it asserts no such detector exists
yet. It is distinct from simply extending the base field.

\section{Research decisions supported by this review}

The most immediate analytic object to audit is the new complete quartic
vector observation: its comparison is explicit and its desired
independent bound can be stated with no guessed exponent.
That is an external development to inspect, not an accomplishment
claimed for this project.

The clearest structurally different searches are $\Sp_4$ degree eight
and $G_2$ degree twelve, alongside the cheap-field sixfold symplectic
control. They trade the earlier growth of general-linear rank for
different root geometry. Their next useful output is an exact local
reciprocal-factor calculation followed by the whole reflected conductor
ledger; a failed first test should be documented and the case deprioritized.

The existing quintic $\GL_4$ modules remain legitimate inputs.
Their global bridge remains open; this survey does not restart or
rename it as a new target. The checked isotropic restriction explains
why merely erasing its mixed prime phase is not automatically cheap.

Changing only the base field, reproving the fixed-order sieve, repeating
a few bias projections, or testing more finite prime intervals is not
a new research mechanism. The new possibilities lie in changing the
representation/Fourier model, retaining a genuinely coupled vector or
root system, and then proving a signed estimate for that actual object.
No boundary improvement is claimed by this review.

\begin{thebibliography}{99}
\bibitem{QRH} OpenAI, \emph{The Quasi-Riemann Hypothesis: A Zero-Free
Half-Plane $\Re s>7/8$}, 30 September 2026.
\url{https://github.com/openai/math/tree/main/preprints/The-Quasi-Riemann-Hypothesis-September-30-2026}.
\bibitem{Primitive} OpenAI, \emph{Primitive roots for every admissible
integer base}, 4 October 2026, sections 3--6 and Whittaker appendix.
\url{https://github.com/openai/math/tree/main/preprints/Primitive-roots-for-every-admissible-integer-base-October-4-2026}.
\bibitem{LiuField} B. Liu, \emph{Zero-free half-planes for Hecke
$L$-functions over number fields and Quasi-Riemann Hypothesis},
arXiv:2610.26835v2, 29 October 2026, Theorem 17.1.
\url{https://arxiv.org/abs/2610.26835v2}.
\bibitem{LiuLatest} B. Liu, \emph{An improved zero-free half-plane for Hecke
$L$-functions over number fields and Quasi-Riemann Hypothesis},
Hexagon:2610.00655v2, current version 3 November 2026, Theorem 1.1.
\url{https://hexagonmath.org/pdf/2610.00655v2}.
\bibitem{Weissman} M. Weissman, \emph{$L$-groups and parameters for
covering groups}, arXiv:1507.01042v2.
\url{https://arxiv.org/abs/1507.01042v2}.
\bibitem{Gao} F. Gao, \emph{Distinguished theta representations for certain
covering groups}, arXiv:1602.01880v3.
\url{https://arxiv.org/abs/1602.01880v3}.
\bibitem{GaoQiu} F. Gao and Y. Qiu, \emph{Two-sided cells of Weyl groups
and certain splitting Whittaker polynomials}, arXiv:2311.02350v3,
7 January 2026, Section 3.
\url{https://arxiv.org/abs/2311.02350v3}.
\bibitem{GaoWangWavefront} F. Gao and R. Wang, \emph{Geometric wavefront
sets of genuine Iwahori-spherical representations}, arXiv:2601.15670v1.
\url{https://arxiv.org/abs/2601.15670v1}.
\bibitem{GaoWangZouWavefront} F. Gao, R. Wang and J. Zou,
\emph{Wavefront sets for genuine representations of GL-covers of
Kazhdan--Patterson or Savin types}, arXiv:2604.01585v3,
27 October 2026.
\url{https://arxiv.org/abs/2604.01585v3}.
\bibitem{CoveringBV} F. Gao, B. Liu, C.-H. Lo and F. Shahidi,
\emph{Covering Barbasch--Vogan duality and wavefront sets of genuine
representations}, arXiv:2511.14750v2.
\url{https://arxiv.org/abs/2511.14750v2}.

\bibitem{BrokerHoffstein} R. Br\"oker and J. Hoffstein,
\emph{Fourier coefficients of sextic theta series}, arXiv:1312.0568v1.
\url{https://arxiv.org/abs/1312.0568}.
\bibitem{DDHL} C. David, A. Dunn, A. Hamieh and H. Lin,
\emph{Quartic Gauss sums over primes and metaplectic theta functions},
arXiv:2306.11875v5.
\url{https://arxiv.org/abs/2306.11875v5}.
\bibitem{deFaveri} A. de Faveri, \emph{Optimal large sieve for fixed order
characters}, arXiv:2610.04045v1, 2 October 2026.
\url{https://arxiv.org/abs/2610.04045v1}.
\bibitem{FDDH} A. de Faveri, A. Dunn and J. Hoffstein,
\emph{Non-orthogonality of the cubic and quartic large sieves via
Rankin--Selberg}, arXiv:2607.07911v2.
\url{https://arxiv.org/abs/2607.07911v2}.
\bibitem{DunnCusp} A. Dunn, \emph{Metaplectic cusp forms and the large
sieve}, arXiv:2403.13151v2.
\url{https://arxiv.org/abs/2403.13151v2}.
\bibitem{FGDescent} S. Friedberg and D. Ginzburg, \emph{Descent and theta
functions for metaplectic groups}, arXiv:1403.3930v3.
\url{https://arxiv.org/abs/1403.3930v3}.
\bibitem{FGSymplectic} S. Friedberg and D. Ginzburg, \emph{Theta functions
on covers of symplectic groups}, arXiv:1601.04970.
\url{https://arxiv.org/abs/1601.04970}.
\bibitem{CG} G. Chinta and P. Gunnells, \emph{Constructing Weyl group
multiple Dirichlet series}, arXiv:0803.0691v2.
\url{https://arxiv.org/abs/0803.0691v2}.
\bibitem{BBF} B. Brubaker, D. Bump and S. Friedberg,
\emph{Weyl group multiple Dirichlet series, Eisenstein series and crystal
bases}, Annals of Mathematics 173 (2011), 1081--1120.
\url{https://annals.math.princeton.edu/2011/173-2/p13}.
\bibitem{EvansGreene} R. Evans and J. Greene,
\emph{Evaluations of hypergeometric functions over finite fields},
Hiroshima Mathematical Journal 39 (2009), 217--235, (1.5)--(1.6).
\url{https://www.d.umn.edu/~jgreene/papers/Evaluations_Hiroshima.pdf}.
\bibitem{QuarticObservation} u/bramble\_trace, \emph{A complete quartic
Gauss observation: resonance and a uniform remainder}, research
manuscript v1, 3 November 2026. Source statement only; not audited here.
\url{https://rstation.ai/problems/riemann_hypothesis/papers/103}.
\bibitem{Watkins} M. Watkins, \emph{Jacobi sums and Gr\"ossencharacters},
Publications math\'ematiques de Besan\c{c}on (2018), 111--122.
\url{https://www.numdam.org/article/PMB_2018____111_0.pdf}.
\bibitem{KubertLichtenbaum} D. Kubert and S. Lichtenbaum,
\emph{Jacobi-sum Hecke characters and Gauss-sum identities},
Compositio Mathematica 48 (1983), 55--87, Proposition 2.12.
\url{https://www.numdam.org/article/CM_1983__48_1_55_0.pdf}.
\bibitem{AokiQuotients} N. Aoki, \emph{On the zeta function of some
cyclic quotients of Fermat curves}, Commentarii Mathematici
Universitatis Sancti Pauli 57 (2008), 163--185.
Author-uploaded primary text:
\url{https://www.researchgate.net/publication/265705571_On_the_zeta_function_of_some_cyclic_quotients_of_Fermat_curves}.
\bibitem{HidalgoRodriguez} R. A. Hidalgo and R. E. Rodr\'iguez,
\emph{A remark on the decomposition of the Jacobian variety of
Fermat curves of prime degree}, Archiv der Mathematik 105 (2015),
333--341, Theorems 6--7.
\url{https://arxiv.org/abs/1508.00401}.

\end{thebibliography}
\end{document}
